\documentclass[11pt]{book}
\usepackage[paperwidth=145mm,paperheight=200mm,top=18mm,bottom=20mm,inner=20mm,outer=16mm]{geometry}
\usepackage{brumfiel}
% figures12.tex -- Chapter 12 figures, built 2026-08-17 from the iPad scans
% (scan14 pp. 8-21 = book pp. 210-226; every figure crop-checked at 150-200 dpi).
%
% Constrained points are DERIVED, never eyeballed:
%   arc midpoints    -> half-angle arithmetic on the central angle
%   angle bisectors  -> half-angle arithmetic
%   crossings        -> (intersection of A--B and C--D)
%   tangency         -> radius-perpendicular construction
%   external points  -> derived as the crossing of the two secants
% Hypotheses are re-asserted numerically in tools/constraints/ch12.py.
%
% 2026-08-17 figure-repair pass (feedback/FIX-SPEC.md).  Figs 12-28, 12-29 and
% 12-30 now carry the BOOK's schematic proportions, NOT the lengths quoted in
% Exercises 15-18: drawing those lengths to scale handed the student the answer.
% Only the structural hypotheses (tangency, collinearity, betweenness, which
% points lie on the circle) are enforced there now.
%
% Macro names are chapter-prefixed: \FIGXII<NUMBERWORD>.
%
% Line treatment follows the book's own convention: solid where the book
% prints solid (given), dashed where the book prints dashed (aux).
%
% Labelling rule used throughout: a point ON the circle is labelled RADIALLY
% OUTWARD, which is always clear of the chords and radii that meet it; interior
% points are labelled along the bisector of the widest free wedge at that point.

\usetikzlibrary{decorations.pathreplacing,arrows.meta}

% Arc-measure marker used in Figs 12-14 to 12-17, and the angle arcs of
% 12-21/12-22.  The book strikes all of these MUCH thinner than the figure
% proper -- the sector's radii and arc read bold against them.
\tikzset{
  XIIhook/.style={fig, line width=0.45pt},
}

% A curly BRACE struck along a circular arc, which is how the book marks an
% arc's measure (Figs 12-14, 12-15, 12-16): a thin curve concentric with the
% arc, an inward hook at each end, and a nib at the middle pointing outward at
% the label.  Before this pass these were drawn as plain arcs with straight
% end ticks (or with nothing at all), which is what Caleb flagged.
%   #1 centre   #2 radius   #3 start angle   #4 end angle   #5 hook/nib depth
% The nib offset is signed off (#4-#3), so the argument order may run either way.
\newcommand{\XIIbrace}[5]{%
  \pgfmathsetmacro{\XIIbm}{0.5*((#3)+(#4))}%
  \pgfmathsetmacro{\XIIbd}{0.05*((#4)-(#3))}%
  \draw[XIIhook, line join=round]
      ($(#1)+({#3}:{(#2)-(#5)})$)
   -- ($(#1)+({#3}:{#2})$)
      arc[start angle={#3}, end angle={\XIIbm-\XIIbd}, radius={#2}]
   -- ($(#1)+({\XIIbm}:{(#2)+(#5)})$)
   -- ($(#1)+({\XIIbm+\XIIbd}:{#2})$)
      arc[start angle={\XIIbm+\XIIbd}, end angle={#4}, radius={#2}]
   -- ($(#1)+({#4}:{(#2)-(#5)})$);%
}

% Primes and degree marks carry a hair kern (1.3mu, about 0.65pt) so that the
% label audit tokenises "A'" and "n^\circ" as single symbols instead of two
% touching glyphs. The shift is well under a point and does not read on the page.
\newcommand{\XIIp}{\mkern1.3mu'}
\newcommand{\XIId}{\mkern1.3mu^\circ}

%----------------------------------------------------------------- Fig 12-1
% Diameter CB through O; A on the lower-left arc; radius OA.  The arc ABC is
% the arc exterior to the central angle AOC.
\newcommand{\FIGXIIONE}{\begin{bkfigure}[\linewidth]
\begin{tikzpicture}[scale=1.300]
  \coordinate (O) at (0,0);
  \def\R{1.25}
  \coordinate (C) at ($(O)+(180:\R)$);
  \coordinate (B) at ($(O)+(0:\R)$);
  \coordinate (A) at ($(O)+(227:\R)$);
  \draw[given] (O) circle[radius=\R];
  \draw[given] (C) -- (B);
  \draw[given] (O) -- (A);
  \node[vlab] at ($(O)+(180:\R+0.28)$) {$C$};
  \node[vlab] at ($(O)+(0:\R+0.28)$) {$B$};
  \node[vlab] at ($(O)+(227:\R+0.28)$) {$A$};
  \node[vlab] at ($(O)+(80:0.32)$) {$O$};
\end{tikzpicture}
\figcap{12--1}
\end{bkfigure}}

%---------------------------------------------------------------- Fig 12-2
% Rays OA, OC, OB with C interior to angle AOB, so angle AOB > angle COB.
% Split out of the old \FIGXIITWOTHREE for the 2026-08-17 pass: 12-2 belongs to
% Exercise 1 of Group 12-1 and 12-3 to Exercise 4, so they sit inline after
% different exercises and can no longer share one captioned block.
\newcommand{\FIGXIITWO}{\begin{bkfigure}[\linewidth]
\begin{tikzpicture}[scale=1.117]
  \coordinate (O) at (0,0);
  \def\R{1.45}
  \coordinate (A) at ($(O)+(78:\R)$);
  \coordinate (C) at ($(O)+(47:\R)$);
  \coordinate (B) at ($(O)+(0:\R)$);
  \draw[given] (O) circle[radius=\R];
  \draw[given] (O) -- (A);
  \draw[given] (O) -- (C);
  \draw[given] (O) -- (B);
  \node[vlab] at ($(O)+(78:\R+0.28)$) {$A$};
  \node[vlab] at ($(O)+(47:\R+0.28)$) {$C$};
  \node[vlab] at ($(O)+(0:\R+0.28)$) {$B$};
  \node[vlab] at ($(O)+(225:0.34)$) {$O$};
\end{tikzpicture}
\figcap{12--2}
\end{bkfigure}}

%---------------------------------------------------------------- Fig 12-3
% Chord EF; P and P' are where rays OQ, OQ' cut EF -- both DERIVED as
% intersections, never placed by hand.
\newcommand{\FIGXIITHREE}{\begin{bkfigure}[\linewidth]
\begin{tikzpicture}[scale=1.117]
  \coordinate (O) at (0,0);
  \def\R{1.90}
  \coordinate (E) at ($(O)+(75:\R)$);
  \coordinate (F) at ($(O)+(0:\R)$);
  \coordinate (Q) at ($(O)+(56:\R)$);
  % Q' at 27 rather than the plate's 22: the wedge Q'OF has to hold the label
  % P', and 22 degrees on this radius leaves it touching the ray.
  \coordinate (Qp) at ($(O)+(27:\R)$);
  \coordinate (P)  at (intersection of O--Q and E--F);
  \coordinate (Pp) at (intersection of O--Qp and E--F);
  \draw[given] (O) circle[radius=\R];
  \draw[given] (O) -- (E);
  \draw[given] (O) -- (F);
  \draw[given] (E) -- (F);
  \draw[given] (O) -- (Q);
  \draw[given] (O) -- (Qp);
  \node[vlab] at ($(O)+(75:\R+0.28)$) {$E$};
  \node[vlab] at ($(O)+(56:\R+0.28)$) {$Q$};
  \node[vlab] at ($(O)+(27:\R+0.30)$) {$Q\XIIp$};
  \node[vlab] at ($(O)+(0:\R+0.28)$) {$F$};
  \node[vlab] at ($(O)+(238:0.34)$) {$O$};
  \node[vlab] at ($(P)+(271.75:0.36)$) {$P$};
  % P' sits on the bisector of the wedge Q'OF, the one direction that clears
  % both the ray OQ' and the chord EF, which cross at P' itself.
  \node[vlab] at ($(Pp)+(250:0.46)$) {$P\XIIp$};
\end{tikzpicture}
\figcap{12--3}
\end{bkfigure}}

%----------------------------------------------------------------- Fig 12-4
% C is the midpoint of the arc EXTERIOR to central angle AOB: the exterior arc
% runs A(118) counterclockwise to B(363), so C sits at (118+363)/2 = 240.5.
\newcommand{\FIGXIIFOUR}{\begin{bkfigure}[\linewidth]
\begin{tikzpicture}[scale=1.300]
  \coordinate (O) at (0,0);
  \def\R{1.25}
  \coordinate (A) at ($(O)+(118:\R)$);
  \coordinate (B) at ($(O)+(3:\R)$);
  \coordinate (C) at ($(O)+(240.5:\R)$);
  \draw[given] (O) circle[radius=\R];
  \draw[given] (O) -- (A);
  \draw[given] (O) -- (B);
  \draw[aux]   (O) -- (C);
  \node[vlab] at ($(O)+(118:\R+0.28)$) {$A$};
  \node[vlab] at ($(O)+(3:\R+0.28)$) {$B$};
  \node[vlab] at ($(O)+(240.5:\R+0.28)$) {$C$};
  \node[vlab] at ($(O)+(60:0.34)$) {$O$};
  \node[slab] at ($(O)+(240.5:\R+0.78)$) {Midpoint};
\end{tikzpicture}
\figcap{12--4}
\end{bkfigure}}

%------------------------------------------------------- Figs 12-5 and 12-6
% The book prints these two side by side on p.211, so they stay one block.
%
% 12-5, rebuilt 2026-08-17 from the book plate (and Caleb's photo of it):
%   * the book draws the two RADII OA and OB, meeting at a shallow 170-degree
%     bend at O -- not the single chord AB the old version drew, which missed
%     the centre by 0.44 r and left O stranded;
%   * the arc therefore spans A(145) to B(-25), 170 degrees, which is what
%     gives the nine labels room.  They were crowded before because the arc
%     was only 146 degrees wide.
%   * C = midpoint of arc AB = 60.  Every later point is the midpoint of the
%     arc it bisects, so all nine positions are arithmetic:
%       D = mid(A,C) = 102.5    E = mid(C,B) = 17.5
%       F = mid(A,D) = 123.75   G = mid(D,C) = 81.25
%       H = mid(C,E) = 38.75    I = mid(E,B) = -3.75
%   * the eight-side path is struck heavier, as the book strikes it, and each
%     of its interior vertices carries a short radial nick.  Caleb's complaint
%     was that the letters read as a caption strip because nothing on the arc
%     marked where the points actually were; the nicks fix that.
% 12-6: same arc, showing an inscribed path against a circumscribed path whose
% corners are the intersections of the tangents at A, C and at C, B --
% distance R/cos(42.5) along the bisecting directions 102.5 and 17.5.
\newcommand{\FIGXIIFIVESIX}{\begin{bkfigure}[\linewidth]
\begin{minipage}[b]{0.54\linewidth}\centering
\begin{tikzpicture}[scale=1.130]
  \coordinate (O) at (0,0);
  \def\R{1.5}
  \coordinate (A) at ($(O)+(145:\R)$);
  \coordinate (B) at ($(O)+(-25:\R)$);
  \coordinate (C) at ($(O)+(60:\R)$);
  \coordinate (D) at ($(O)+(102.5:\R)$);
  \coordinate (E) at ($(O)+(17.5:\R)$);
  \coordinate (F) at ($(O)+(123.75:\R)$);
  \coordinate (G) at ($(O)+(81.25:\R)$);
  \coordinate (H) at ($(O)+(38.75:\R)$);
  \coordinate (I) at ($(O)+(-3.75:\R)$);
  \draw[given] (O) circle[radius=\R];
  \draw[given] (O) -- (A);
  \draw[given] (O) -- (B);
  % the eight-side path is struck HEAVIER in the book, not dashed, so that it
  % reads apart from the dashed two-side path it refines
  \draw[fig, line width=1.5pt]
        (A) -- (F) -- (D) -- (G) -- (C) -- (H) -- (E) -- (I) -- (B);
  % short radial nicks at the interior vertices, so each letter labels a place
  % on the arc and not a slot in a row
  \foreach \a in {123.75, 102.5, 81.25, 38.75, 17.5, -3.75}
    {\draw[fig, line width=0.5pt] ($(O)+(\a:\R)$) -- ($(O)+(\a:\R+0.09)$);}
  \draw[aux] (A) -- (C) -- (B);
  \foreach \p/\a in {A/145, B/-25, C/60, D/102.5, E/17.5,
                     F/123.75, G/81.25, H/38.75, I/-3.75}
    {\node[vlab] at ($(O)+(\a:\R+0.38)$) {$\p$};}
  \node[vlab] at ($(O)+(238:0.34)$) {$O$};
\end{tikzpicture}
\figcap{12--5}
\end{minipage}\hfill
\begin{minipage}[b]{0.40\linewidth}\centering
\begin{tikzpicture}[scale=1.060]
  \coordinate (O) at (0,0);
  \def\R{1.3}
  \coordinate (A) at ($(O)+(145:\R)$);
  \coordinate (B) at ($(O)+(-25:\R)$);
  \coordinate (C) at ($(O)+(60:\R)$);
  % tangent-intersection corners of the circumscribed path
  \pgfmathsetmacro{\Rt}{\R/cos(42.5)}
  \coordinate (T) at ($(O)+(102.5:\Rt)$);
  \coordinate (U) at ($(O)+(17.5:\Rt)$);
  \draw[given] (O) circle[radius=\R];
  \draw[given] (O) -- (A);
  \draw[given] (O) -- (B);
  \draw[aux] (A) -- (C) -- (B);
  \draw[aux] (A) -- (T) -- (C) -- (U) -- (B);
\end{tikzpicture}
\figcap{12--6}
\end{minipage}
\end{bkfigure}}

%---------------------------------------------------------------- Fig 12-7
% Regular hexagon, A(120), C(60), B(0); radii OA, OC, OB drawn.
% Split out of the old \FIGXIISEVENEIGHT: 12-7 belongs to Exercise 1 of Group
% 12-2 and 12-8 to Exercise 3, so each now stands inline after its own item.
\newcommand{\FIGXIISEVEN}{\begin{bkfigure}[\linewidth]
\begin{tikzpicture}[scale=1.196]
  \coordinate (O) at (0,0);
  \def\R{1.25}
  \coordinate (A) at ($(O)+(120:\R)$);
  \coordinate (C) at ($(O)+(60:\R)$);
  \coordinate (B) at ($(O)+(0:\R)$);
  \draw[given] (O) circle[radius=\R];
  \draw[given] ($(O)+(0:\R)$) -- ($(O)+(60:\R)$) -- ($(O)+(120:\R)$)
            -- ($(O)+(180:\R)$) -- ($(O)+(240:\R)$) -- ($(O)+(300:\R)$) -- cycle;
  \draw[given] (O) -- (A);
  \draw[given] (O) -- (C);
  \draw[given] (O) -- (B);
  \node[vlab] at ($(O)+(120:\R+0.30)$) {$A$};
  \node[vlab] at ($(O)+(60:\R+0.30)$) {$C$};
  \node[vlab] at ($(O)+(0:\R+0.30)$) {$B$};
  \node[vlab] at ($(O)+(255:0.34)$) {$O$};
\end{tikzpicture}
\figcap{12--7}
\end{bkfigure}}

%---------------------------------------------------------------- Fig 12-8
% Angle AOC is a right angle and angle AOB = angle BOC, so B bisects it.
% Rebuilt 2026-08-17.  Two things were wrong and Caleb caught both:
%   * the OUTER arc ran from OA all the way to OC, CROSSING ray OB -- it marked
%     angle AOC.  The book strikes TWO arcs that abut at OB, one on angle AOB
%     and one on angle BOC, at clearly different radii (the book's differ by
%     about 60%); nothing crosses a ray.  That adjacency is the whole point of
%     the exercise.
%   * the fan was axis-aligned (OA vertical, OC horizontal), so the render read
%     as though angle AOC were the axis frame rather than the right angle the
%     exercise names.  The book tilts the fan; A(78), B(33), C(-12) keeps
%     angle AOC = 90 and angle AOB = angle BOC = 45 exactly, with no ray level.
\newcommand{\FIGXIIEIGHT}{\begin{bkfigure}[\linewidth]
\begin{tikzpicture}[scale=1.196]
  \coordinate (O) at (0,0);
  \def\R{1.25}
  \coordinate (A) at ($(O)+(78:\R)$);
  \coordinate (B) at ($(O)+(33:\R)$);
  \coordinate (C) at ($(O)+(-12:\R)$);
  \draw[given] (O) circle[radius=\R];
  \draw[given] (O) -- (A);
  \draw[given] (O) -- (B);
  \draw[given] (O) -- (C);
  % angle AOB, inner radius; angle BOC, outer radius.  They meet on ray OB.
  \draw[XIIhook] ($(O)+(33:0.36)$) arc[start angle=33, end angle=78, radius=0.36];
  \draw[XIIhook] ($(O)+(-12:0.60)$) arc[start angle=-12, end angle=33, radius=0.60];
  \node[vlab] at ($(O)+(78:\R+0.28)$) {$A$};
  \node[vlab] at ($(O)+(33:\R+0.28)$) {$B$};
  \node[vlab] at ($(O)+(-12:\R+0.28)$) {$C$};
  \node[vlab] at ($(O)+(228:0.28)$) {$O$};
\end{tikzpicture}
\figcap{12--8}
\end{bkfigure}}

%----------------------------------------------------- Figs 12-9 and 12-10
% 12-9: an IRREGULAR inscribed polygonal path in arc AB (Theorem 12-2 drops the
%   regularity requirement), so the vertex spacing is deliberately uneven.
%   Rebuilt 2026-08-17: the path used to have SEVEN short sides, whose sag off
%   the circle was 0.017 r -- under a point on the page.  The dashes therefore
%   printed on top of the circle line and the whole thing read as a ragged
%   doubled edge, which is what Caleb saw.  The book draws four long sides,
%   whose sag is about 0.07 r, so the dashed path is plainly INSIDE the circle
%   with a visible gap.  Spacing 40/34/48/38 degrees keeps it irregular.
% 12-10: adjacent arcs AB and BC, sharing only B.  The emphasis is now a
%   heavier DASH-BROKEN arc struck just outside the circle, which is how the
%   book marks it; simply thickening the circle's own stroke (the old \key arc)
%   was invisible at 1.0pt against a 0.9pt circle.  The fan is tilted too --
%   the book sets no radius on an axis.
\newcommand{\FIGXIININETEN}{\begin{bkfigure}[\linewidth]
\begin{minipage}[b]{0.46\linewidth}\centering
\begin{tikzpicture}[scale=1.198]
  \coordinate (O) at (0,0);
  \def\R{1.25}
  \coordinate (A) at ($(O)+(152:\R)$);
  \coordinate (B) at ($(O)+(-8:\R)$);
  \draw[given] (O) circle[radius=\R];
  \draw[given] (O) -- (A);
  \draw[given] (O) -- (B);
  \draw[aux] (A)
    -- ($(O)+(112:\R)$) -- ($(O)+(78:\R)$) -- ($(O)+(30:\R)$) -- (B);
  \node[vlab] at ($(O)+(152:\R+0.28)$) {$A$};
  \node[vlab] at ($(O)+(-8:\R+0.28)$) {$B$};
  \node[vlab] at ($(O)+(248:0.30)$) {$O$};
\end{tikzpicture}
\figcap{12--9}
\end{minipage}\hfill
\begin{minipage}[b]{0.46\linewidth}\centering
\begin{tikzpicture}[scale=1.198]
  \coordinate (O) at (0,0);
  \def\R{1.25}
  \coordinate (A) at ($(O)+(82:\R)$);
  \coordinate (B) at ($(O)+(37:\R)$);
  \coordinate (C) at ($(O)+(-8:\R)$);
  \draw[given] (O) circle[radius=\R];
  \draw[given] (O) -- (A);
  \draw[given] (O) -- (B);
  \draw[given] (O) -- (C);
  \draw[fig, line width=1.5pt, dash pattern=on 3.6pt off 1.5pt]
        ($(O)+(-8:\R+0.05)$) arc[start angle=-8, end angle=82, radius=\R+0.05];
  \node[vlab] at ($(O)+(82:\R+0.30)$) {$A$};
  \node[vlab] at ($(O)+(37:\R+0.30)$) {$B$};
  \node[vlab] at ($(O)+(-8:\R+0.30)$) {$C$};
  \node[vlab] at ($(O)+(222:0.30)$) {$O$};
\end{tikzpicture}
\figcap{12--10}
\end{minipage}
\end{bkfigure}}

%----------------------------------------------------- Figs 12-11 and 12-12
% 12-11: C interior to angle AOB, so angle AOB (93) > angle COB (80).
%   Directions measured off the plate: A 104, C 90, B 10 (the book sets A and
%   C only about 14 degrees apart, which is why their labels sit side by side).
% 12-12: corresponding central angles congruent on both circles -- the same
% four directions 148, 90, 48, 0 are used at each centre, which is exactly the
% hypothesis of Theorem 12-7.
\newcommand{\FIGXIIELEVENTWELVE}{\begin{bkfigure}[\linewidth]
\begin{minipage}[b]{0.32\linewidth}\centering
\begin{tikzpicture}[scale=0.995]
  \coordinate (O) at (0,0);
  \def\R{1.25}
  \coordinate (A) at ($(O)+(104:\R)$);
  \coordinate (C) at ($(O)+(90:\R)$);
  \coordinate (B) at ($(O)+(10:\R)$);
  \draw[given] (O) circle[radius=\R];
  \draw[given] (O) -- (A);
  \draw[given] (O) -- (C);
  \draw[given] (O) -- (B);
  \node[vlab] at ($(O)+(112:\R+0.32)$) {$A$};
  \node[vlab] at ($(O)+(86:\R+0.32)$) {$C$};
  \node[vlab] at ($(O)+(10:\R+0.30)$) {$B$};
  \node[vlab] at ($(O)+(250:0.34)$) {$O$};
\end{tikzpicture}
\figcap{12--11}
\end{minipage}\hfill
\begin{minipage}[b]{0.64\linewidth}\centering
\begin{tikzpicture}[scale=0.922]
  \coordinate (Op) at (0,0);
  \def\Rp{1.35}
  \coordinate (Ap) at ($(Op)+(148:\Rp)$);
  \coordinate (Pp) at ($(Op)+(90:\Rp)$);
  \coordinate (Qp) at ($(Op)+(48:\Rp)$);
  \coordinate (Bp) at ($(Op)+(0:\Rp)$);
  \draw[given] (Op) circle[radius=\Rp];
  \draw[given] (Op) -- (Ap);
  \draw[given] (Op) -- (Bp);
  \draw[aux] (Op) -- (Pp);
  \draw[aux] (Op) -- (Qp);
  \draw[aux] (Ap) -- (Pp) -- (Qp) -- (Bp);
  \node[vlab] at ($(Op)+(152:\Rp+0.40)$) {$A\XIIp$};
  \node[vlab] at ($(Op)+(90:\Rp+0.34)$) {$P\XIIp$};
  \node[vlab] at ($(Op)+(44:\Rp+0.40)$) {$Q\XIIp$};
  \node[vlab] at ($(Op)+(0:\Rp+0.32)$) {$B\XIIp$};
  \node[vlab] at ($(Op)+(250:0.48)$) {$O\XIIp$};
  \node[slab] at ($(Op)+(0:0.68)+(0,-0.40)$) {$r\XIIp$};
  %
  \coordinate (Oo) at (3.75,0);
  \def\Rr{0.85}
  \coordinate (Aa) at ($(Oo)+(148:\Rr)$);
  \coordinate (Pq) at ($(Oo)+(90:\Rr)$);
  \coordinate (Qq) at ($(Oo)+(48:\Rr)$);
  \coordinate (Bb) at ($(Oo)+(0:\Rr)$);
  \draw[given] (Oo) circle[radius=\Rr];
  \draw[given] (Oo) -- (Aa);
  \draw[given] (Oo) -- (Bb);
  \draw[aux] (Oo) -- (Pq);
  \draw[aux] (Oo) -- (Qq);
  \draw[aux] (Aa) -- (Pq) -- (Qq) -- (Bb);
  \node[vlab] at ($(Oo)+(148:\Rr+0.30)$) {$A$};
  \node[vlab] at ($(Oo)+(90:\Rr+0.30)$) {$P$};
  \node[vlab] at ($(Oo)+(48:\Rr+0.30)$) {$Q$};
  \node[vlab] at ($(Oo)+(0:\Rr+0.30)$) {$B$};
  \node[vlab] at ($(Oo)+(250:0.36)$) {$O$};
  \node[slab] at ($(Oo)+(0:0.44)+(0,-0.38)$) {$r$};
\end{tikzpicture}
\figcap{12--12}
\end{minipage}
\end{bkfigure}}

%---------------------------------------------------- Figs 12-13 and 12-14
% 12-13: OC bisects the right angle AOB, OD bisects BOC, OE bisects DOC.
%   A = 0, B = 90, C = 45, D = 67.5, E = 56.25 -- pure halving, no eyeballing.
% 12-14: rebuilt 2026-08-17 from the book plate.  Four faults, all Caleb's:
%   * it was drawn as an axis-aligned 90-degree quarter.  The book draws a
%     roughly 68-degree sector, tilted so it is symmetric about the horizontal
%     and neither bounding radius is level: A(-34), B(+34).
%   * it carried five interior rays running clear THROUGH the outer marker, so
%     the crossings read as tick marks on it and the outer region looked
%     hatched.  The book has exactly four, overshooting the arc only slightly
%     and stopping well short of the marker.
%   * the interior rays are the same subdivisions 12-13 constructs, read as
%     fractions of the sector measured up from OA: 1/2, 5/8, 3/4, all dashed,
%     plus the SOLID ray at 0.662 -- "a little greater than 5/8", exactly what
%     the text claims of L.  That is why the book's third ray sits a hair below
%     the solid one.
%   * the l marker was a plain arc spanning the WHOLE sector with no hooks.  It
%     is a brace, it carries end hooks and a middle nib, and it spans only the
%     sub-arc AP whose length is l.
\newcommand{\FIGXIITHIRTEENFOURTEEN}{\begin{bkfigure}[\linewidth]
\begin{minipage}[b]{0.48\linewidth}\centering
\begin{tikzpicture}[scale=1.194]
  \coordinate (O) at (0,0);
  \def\R{1.6}
  \coordinate (A) at ($(O)+(0:\R)$);
  \coordinate (B) at ($(O)+(90:\R)$);
  \coordinate (C) at ($(O)+(45:\R)$);
  \coordinate (D) at ($(O)+(67.5:\R)$);
  \coordinate (E) at ($(O)+(56.25:\R)$);
  \draw[given] (O) circle[radius=\R];
  \draw[given] (O) -- (B);
  \draw[given] (O) -- (A);
  \draw[aux] (O) -- (D);
  \draw[aux] (O) -- (E);
  \draw[aux] (O) -- (C);
  \node[vlab] at ($(O)+(90:\R+0.28)$) {$B$};
  \node[vlab] at ($(O)+(67.5:\R+0.28)$) {$D$};
  \node[vlab] at ($(O)+(56.25:\R+0.28)$) {$E$};
  \node[vlab] at ($(O)+(45:\R+0.28)$) {$C$};
  \node[vlab] at ($(O)+(0:\R+0.28)$) {$A$};
  \node[vlab] at ($(O)+(250:0.34)$) {$O$};
  \node[slab] at ($(O)+(90:0.80)+(-0.28,0)$) {$r$};
\end{tikzpicture}
\figcap{12--13}
\end{minipage}\hfill
\begin{minipage}[b]{0.44\linewidth}\centering
\begin{tikzpicture}[scale=1.099]
  \coordinate (O) at (0,0);
  \def\R{1.7}
  % the sector opens 68 degrees, symmetric about the horizontal
  \def\Alo{-34}\def\Bhi{34}
  \coordinate (A) at ($(O)+(\Alo:\R)$);
  \coordinate (B) at ($(O)+(\Bhi:\R)$);
  % interior rays as fractions of the opening measured up from OA
  \pgfmathsetmacro{\rhalf}{\Alo+0.500*(\Bhi-\Alo)}
  \pgfmathsetmacro{\rfive}{\Alo+0.625*(\Bhi-\Alo)}
  \pgfmathsetmacro{\rP}{\Alo+0.662*(\Bhi-\Alo)}
  \pgfmathsetmacro{\rthree}{\Alo+0.750*(\Bhi-\Alo)}
  \coordinate (P) at ($(O)+(\rP:\R)$);
  \draw[given] (O) -- (B);
  \draw[given] (O) -- (A);
  \draw[given] ($(O)+(\Alo:\R)$)
      arc[start angle=\Alo, end angle=\Bhi, radius=\R];
  \foreach \a in {\rhalf, \rfive, \rthree}
    {\draw[aux] (O) -- ($(O)+(\a:\R+0.15)$);}
  \draw[given] (O) -- ($(O)+(\rP:\R+0.15)$);
  % brace marking the arc AP, of length l: hooks at both ends, nib at the
  % middle pointing out at the label, clear of the rays' overshoot
  \XIIbrace{O}{\R+0.36}{\rP}{\Alo}{0.14}
  \node[vlab] at ($(O)+(\Bhi:\R)+(0.06,0.28)$) {$B$};
  \node[vlab] at ($(O)+(\Alo:\R)+(0.06,-0.28)$) {$A$};
  \node[vlab] at ($(O)+(196:0.30)$) {$O$};
  % r sits INSIDE the sector, against OB, as the book sets it -- the render had
  % it outside, on the far side of the radius.  It is centred in the 17-degree
  % wedge between OB and the 3/4 ray, far enough out that the wedge is wide
  % enough to hold it clear of both.
  \node[slab] at ($(O)+(25.5:1.35)$) {$r$};
  \pgfmathsetmacro{\lmid}{0.5*(\rP+\Alo)}
  \node[slab] at ($(O)+(\lmid:\R+0.74)$) {$l$};
\end{tikzpicture}
\figcap{12--14}
\end{minipage}
\end{bkfigure}}

%---------------------------------------------------------------- Fig 12-15
% Successive 1-degree arcs counted from B toward A, exaggerated for legibility.
% A(48) lies strictly between the n-degree ray (32) and the (n+1)-degree ray
% (57), which is exactly the inequality the theorem asserts.
%
% Rebuilt 2026-08-17.  The two arc-measure markers were plain concentric arcs
% -- the n-degree one with straight end ticks, the (n+1)-degree one with
% nothing at all and its label floating free.  The book draws both as BRACES:
% inward hooks at each end and a nib at the middle pointing straight at the
% label.  The 1-degree marker was a bare line with a crossbar dropped far below
% the figure; the book tucks a short headed arrow into the 1-degree wedge with
% the label immediately under OB.
\newcommand{\FIGXIIFIFTEEN}{\begin{bkfigure}[\linewidth]
\begin{tikzpicture}[scale=1.300]
  \coordinate (O) at (0,0);
  \def\R{1.95}
  \coordinate (B) at ($(O)+(0:\R)$);
  \coordinate (A) at ($(O)+(48:\R)$);
  \draw[given] (O) -- (B);
  \draw[given] (O) -- (A);
  \draw[given] ($(O)+(0:\R)$) arc[start angle=0, end angle=48, radius=\R];
  % each dashed ray runs out only as far as the marker it serves: the 32-degree
  % ray closes the n-degree brace, the 57-degree ray closes the (n+1)-degree
  % brace, and the 10-degree ray merely clears the arc
  \foreach \a/\len in {10/0.10, 32/0.24, 57/0.92}
    {\draw[aux] (O) -- ($(O)+(\a:\R+\len)$);}
  \XIIbrace{O}{\R+0.20}{0}{32}{0.12}
  \XIIbrace{O}{\R+0.80}{0}{57}{0.14}
  % the 1-degree wedge: one short arrow rising from under OB, its head landing
  % inside the wedge between OB and the 10-degree ray
  \draw[XIIhook, -{Stealth[length=4pt,width=3pt]}] (1.28,-0.34) -- (1.28,0.15);
  \node[vlab] at ($(O)+(200:0.34)$) {$O$};
  \node[vlab] at ($(O)+(0:\R)+(0.16,-0.26)$) {$B$};
  \node[vlab] at ($(O)+(48:\R+0.34)$) {$A$};
  \node[slab] at ($(O)+(32:1.45)+(-58:0.32)$) {$r$};
  \node[slab, anchor=west] at ($(O)+(16:\R+0.38)$) {$n\XIId$};
  \node[slab, anchor=west] at ($(O)+(28.5:\R+1.00)$) {$(n+1)\XIId$};
  \node[slab] at (1.28,-0.60) {$1\XIId$};
\end{tikzpicture}
\figcap{12--15}
\end{bkfigure}}

%---------------------------------------------------------------- Fig 12-16
% The primed twin of Fig 12-15: angle A'O'B' congruent to angle AOB, counted
% the same way, so A' falls on the same (n+1)st arc.
% 2026-08-17: the radius was labelled a bare r here.  EVERY letter in this
% figure is primed in the book -- it is the primed twin -- so it is r', and
% Caleb flagged the bare r by name.  Markers rebuilt as braces, as in 12-15.
\newcommand{\FIGXIISIXTEEN}{\begin{bkfigure}[\linewidth]
\begin{tikzpicture}[scale=1.300]
  \coordinate (O) at (0,0);
  \def\R{1.95}
  \coordinate (B) at ($(O)+(0:\R)$);
  \coordinate (A) at ($(O)+(48:\R)$);
  \draw[given] (O) -- (B);
  \draw[given] (O) -- (A);
  \draw[given] ($(O)+(0:\R)$) arc[start angle=0, end angle=48, radius=\R];
  \foreach \a/\len in {10/0.10, 32/0.24, 57/0.92}
    {\draw[aux] (O) -- ($(O)+(\a:\R+\len)$);}
  \XIIbrace{O}{\R+0.20}{0}{32}{0.12}
  \XIIbrace{O}{\R+0.80}{0}{57}{0.14}
  \draw[XIIhook, -{Stealth[length=4pt,width=3pt]}] (1.28,-0.34) -- (1.28,0.15);
  \node[vlab] at ($(O)+(200:0.36)$) {$O\XIIp$};
  \node[vlab] at ($(O)+(0:\R)+(0.20,-0.26)$) {$B\XIIp$};
  \node[vlab] at ($(O)+(48:\R+0.34)$) {$A\XIIp$};
  \node[slab] at ($(O)+(32:1.45)+(-58:0.34)$) {$r\XIIp$};
  \node[slab, anchor=west] at ($(O)+(16:\R+0.38)$) {$n\XIId$};
  \node[slab, anchor=west] at ($(O)+(28.5:\R+1.00)$) {$(n+1)\XIId$};
  \node[slab] at (1.28,-0.60) {$1\XIId$};
\end{tikzpicture}
\figcap{12--16}
\end{bkfigure}}

%--------------------------------------------------- Figs 12-17 and 12-18
% 12-17: left, the arc subtended by a central angle of A degrees; right, the
% arc EXTERIOR to a congruent central angle, of 360 - A degrees.  The two
% central angles are drawn with the same opening (75 degrees).
% 12-18: one radian -- the central angle whose subtended arc equals the radius,
% so the opening is exactly 180/pi = 57.2958 degrees.
\newcommand{\FIGXIISEVENTEENEIGHTEEN}{\begin{bkfigure}[\linewidth]
\begin{minipage}[b]{0.62\linewidth}\centering
\begin{tikzpicture}[scale=0.76]
  % left circle: interior arc
  \coordinate (P) at (0,0);
  % 1.15, not 0.95: the interior A-degree numeral has to sit inside a 75-degree
  % wedge, and on the smaller radius its box overhangs the upper ray.
  \def\Ra{1.15}
  \draw[given] (P) circle[radius=\Ra];
  \draw[given] (P) -- ($(P)+(37.5:\Ra)$);
  \draw[given] (P) -- ($(P)+(-37.5:\Ra)$);
  \draw[given] ($(P)+(-37.5:0.26)$) arc[start angle=-37.5, end angle=37.5, radius=0.20];
  \node[slab] at ($(P)+(0:0.70)$) {$A\XIId$};
  \draw[XIIhook] ($(P)+(-37.5:\Ra+0.22)$)
      arc[start angle=-37.5, end angle=37.5, radius=\Ra+0.22];
  \draw[XIIhook] ($(P)+(37.5:\Ra+0.22)$) -- ($(P)+(37.5:\Ra+0.10)$);
  \draw[XIIhook] ($(P)+(-37.5:\Ra+0.22)$) -- ($(P)+(-37.5:\Ra+0.10)$);
  \node[slab] at ($(P)+(0:\Ra+0.58)$) {$A\XIId$};
  % right circle: exterior arc
  \coordinate (Q) at (3.95,0);
  \def\Rb{1.2}
  \draw[given] (Q) circle[radius=\Rb];
  \draw[given] (Q) -- ($(Q)+(142.5:\Rb)$);
  \draw[given] (Q) -- ($(Q)+(217.5:\Rb)$);
  \draw[given] ($(Q)+(142.5:0.28)$) arc[start angle=142.5, end angle=217.5, radius=0.24];
  \node[slab] at ($(Q)+(180:0.74)$) {$A\XIId$};
  \draw[XIIhook] ($(Q)+(142.5:\Rb+0.24)$)
      arc[start angle=142.5, end angle=-142.5, radius=\Rb+0.24];
  \draw[XIIhook] ($(Q)+(142.5:\Rb+0.24)$) -- ($(Q)+(142.5:\Rb+0.12)$);
  \draw[XIIhook] ($(Q)+(217.5:\Rb+0.24)$) -- ($(Q)+(217.5:\Rb+0.12)$);
  \node[slab, right] at ($(Q)+(0:\Rb+0.34)$) {$(360\XIId-A\XIId\mkern1.3mu)$};
\end{tikzpicture}
\figcap{12--17}
\end{minipage}\hfill
\begin{minipage}[b]{0.34\linewidth}\centering
\begin{tikzpicture}[scale=0.86]
  \coordinate (O) at (0,0);
  \def\R{1.3}
  \coordinate (U) at ($(O)+(50:\R)$);
  \coordinate (V) at ($(O)+(-7.2958:\R)$);
  \draw[given] (O) circle[radius=\R];
  \draw[given] (O) -- (U);
  \draw[given] (O) -- (V);
  \draw[given] ($(O)+(-7.2958:0.44)$)
      arc[start angle=-7.2958, end angle=50, radius=0.38];
  \node[slab] at ($(O)+(21.35:0.82)$) {$1$};
  \node[slab] at ($(O)!0.58!(U)+(-0.30,0.18)$) {$r$};
  \node[slab] at ($(O)+(21.35:\R+0.28)$) {$r$};
\end{tikzpicture}
\figcap{12--18}
\end{minipage}
\end{bkfigure}}

%--------------------------------------------------- Figs 12-19 and 12-20
% 12-19: two inscribed angles (vertex ON the circle) and one that is not
% (vertex OUTSIDE the circle) -- the distinction the definition turns on.
% 12-20: BC a diameter, so the inscribed angle B stands on the arc AC; OA is
% the radius drawn for the lemma's isosceles triangle OAB.
\newcommand{\FIGXIININETEENTWENTY}{\begin{bkfigure}[\linewidth]
\begin{minipage}[b]{0.66\linewidth}\centering
\begin{tikzpicture}[scale=0.74]
  \def\Rc{0.95}
  % (a) inscribed
  \coordinate (Oa) at (0,0);
  \coordinate (Ba) at ($(Oa)+(168:\Rc)$);
  \coordinate (Aa) at ($(Oa)+(12:\Rc)$);
  \coordinate (Ca) at ($(Oa)+(-55:\Rc)$);
  \draw[given] (Oa) circle[radius=\Rc];
  \draw[given] (Ba) -- (Aa);
  \draw[given] (Ba) -- (Ca);
  \node[vlab] at ($(Oa)+(168:\Rc+0.34)$) {$B$};
  \node[vlab] at ($(Aa)+(48:0.36)$) {$A$};
  \node[vlab] at ($(Oa)+(-55:\Rc+0.34)$) {$C$};
  \node[slab] at ($(Oa)+(0,-1.85)$) {Inscribed};
  % (b) inscribed
  \coordinate (Ob) at (3.15,0);
  \coordinate (Bb) at ($(Ob)+(232:\Rc)$);
  \coordinate (Ab) at ($(Ob)+(119:\Rc)$);
  \coordinate (Cb) at ($(Ob)+(-43:\Rc)$);
  \draw[given] (Ob) circle[radius=\Rc];
  \draw[given] (Bb) -- (Ab);
  \draw[given] (Bb) -- (Cb);
  \node[vlab] at ($(Ab)+(136:0.50)$) {$A$};
  \node[vlab] at ($(Ob)+(232:\Rc+0.34)$) {$B$};
  \node[vlab] at ($(Ob)+(-43:\Rc+0.34)$) {$C$};
  \node[slab] at ($(Ob)+(0,-1.85)$) {Inscribed};
  % (c) NOT inscribed.  Measured off scan14 p.17: the vertex B sits a whisker
  % OUTSIDE the circle (1.08 r) and the side BA runs AWAY from it, so the whole
  % ray BA misses the circle -- that is what fails the definition.  Side BC, by
  % contrast, cuts straight across.  (Directions 215.6 / 111.8 / -45.3 and the
  % 1.5 r length of BA are read straight off the plate.)
  \coordinate (Oc) at (6.30,0);
  \coordinate (Bc) at ($(Oc)+(215.6:1.08*\Rc)$);
  \coordinate (Ac) at ($(Bc)+(111.8:1.50*\Rc)$);
  \coordinate (Cc) at ($(Oc)+(-45.3:\Rc)$);
  \draw[given] (Oc) circle[radius=\Rc];
  \draw[given] (Ac) -- (Bc);
  \draw[given] (Bc) -- (Cc);
  \node[vlab] at ($(Ac)+(50:0.34)$) {$A$};
  \node[vlab] at ($(Bc)+(205:0.34)$) {$B$};
  \node[vlab] at ($(Cc)+(-40:0.34)$) {$C$};
  \node[slab] at ($(Oc)+(0,-1.85)$) {Not inscribed};
\end{tikzpicture}
\figcap{12--19}
\end{minipage}\hfill
\begin{minipage}[b]{0.30\linewidth}\centering
\begin{tikzpicture}[scale=0.86]
  \coordinate (O) at (0,0);
  \def\R{1.25}
  \coordinate (B) at ($(O)+(180:\R)$);
  \coordinate (C) at ($(O)+(0:\R)$);
  \coordinate (A) at ($(O)+(79:\R)$);
  \draw[given] (O) circle[radius=\R];
  \draw[given] (B) -- (C);
  \draw[given] (B) -- (A);
  \draw[aux]   (A) -- (O);   % the book prints this radius DASHED
  \node[vlab] at ($(O)+(79:\R+0.30)$) {$A$};
  \node[vlab] at ($(O)+(180:\R+0.30)$) {$B$};
  \node[vlab] at ($(O)+(0:\R+0.30)$) {$C$};
  \node[vlab] at ($(O)+(255:0.36)$) {$O$};
\end{tikzpicture}
\figcap{12--20}
\end{minipage}
\end{bkfigure}}

%--------------------------------------------------- Figs 12-21 and 12-22
% 12-21 (Case 1): the centre is INTERIOR to angle ABC.  The diameter BP splits
%   arc AC into AP(80) and PC(33), and angle ABC into x = 40 and y = 16.5.
%   Directions read off scan14 p.18: A 80, C -33.  The two angle arcs at B are
%   in the book and are what tie the letters x, y to the angles they name.
% 12-22 (Case 2): the centre is EXTERIOR to angle ABC -- both rays BA, BC lie
%   on the same side of the diameter BQ.  Measured: A 110, C 36.
%   Rebuilt 2026-08-17; three faults, all Caleb's, all confirmed against the
%   plate at high zoom:
%     * a chord AC was drawn that the book does NOT have, closing a triangle
%       ABC.  Between A and C the book has only the circle's arc.  Removed.
%     * BOTH angle arcs at B were missing.  The book strikes two thin
%       concentric arcs centred at B, the outer springing from A itself and the
%       inner from about 55% of the way along chord BA; each sweeps clockwise
%       ACROSS chord BC and lands on the dashed diameter BQ with a short
%       downward tick.  They are what carry the subtraction the proof needs,
%       angle ABC = angle ABQ - angle CBQ.
%     * the centre carried no dot.  The book marks it, and sets O above-left.
\newcommand{\FIGXIITWENTYONETWENTYTWO}{\begin{bkfigure}[\linewidth]
\begin{minipage}[b]{0.46\linewidth}\centering
\begin{tikzpicture}[scale=1.091]
  \coordinate (O) at (0,0);
  % 1.6: y names a 16.5-degree wedge, and its numeral needs about 1.5 units of
  % run from B before the wedge is tall enough to hold it clear of both sides.
  \def\R{1.6}
  \coordinate (B) at ($(O)+(180:\R)$);
  \coordinate (P) at ($(O)+(0:\R)$);
  \coordinate (A) at ($(O)+(80:\R)$);
  \coordinate (C) at ($(O)+(-33:\R)$);
  \draw[given] (O) circle[radius=\R];
  \draw[given] (B) -- (A);
  \draw[given] (A) -- (C);
  \draw[given] (B) -- (C);
  \draw[aux] (B) -- (P);
  % angle arcs: x = angle ABP (40 deg), y = angle PBC (16.5 deg)
  \draw[XIIhook] ($(B)+(0:0.60)$)     arc[start angle=0,     end angle=40, radius=0.60];
  \draw[XIIhook] ($(B)+(-16.5:0.90)$) arc[start angle=-16.5, end angle=0,  radius=0.90];
  \node[vlab] at ($(O)+(80:\R+0.30)$) {$A$};
  \node[vlab] at ($(O)+(180:\R+0.30)$) {$B$};
  \node[vlab] at ($(O)+(0:\R+0.30)$) {$P$};
  \node[vlab] at ($(O)+(-33:\R+0.30)$) {$C$};
  \node[vlab] at ($(O)+(90:0.30)$) {$O$};
  \node[slab] at ($(B)+(20:1.00)$) {$x$};
  \node[slab] at ($(B)+(-8.25:1.85)$) {$y$};
\end{tikzpicture}
\figcap{12--21}
\end{minipage}\hfill
\begin{minipage}[b]{0.46\linewidth}\centering
\begin{tikzpicture}[scale=1.091]
  \coordinate (O) at (0,0);
  \def\R{1.35}
  \coordinate (B) at ($(O)+(180:\R)$);
  \coordinate (Q) at ($(O)+(0:\R)$);
  \coordinate (A) at ($(O)+(110:\R)$);
  \coordinate (C) at ($(O)+(36:\R)$);
  \draw[given] (O) circle[radius=\R];
  \draw[given] (B) -- (A);
  \draw[given] (B) -- (C);
  \draw[aux] (B) -- (Q);
  % the two angle arcs at B.  Their common span is angle ABQ, so both are
  % DERIVED from the directions B->A and B->Q, never swept by eye; the outer
  % radius is |BA| exactly, so that arc springs from A itself.
  % chord direction and chord length on a circle, in closed form: the ray from
  % the point at angle p to the point at angle q runs at (p+q)/2 - 90, and the
  % chord measures 2 r sin((p-q)/2).  B is at 180, A at 110.
  \pgfmathsetmacro{\bA}{0.5*(180+110)-90}
  \pgfmathsetmacro{\bBA}{2*\R*sin(0.5*(180-110))}
  \foreach \f in {1.0, 0.55}{%
    \pgfmathsetmacro{\rr}{\f*\bBA}
    \draw[XIIhook] ($(B)+(\bA:\rr)$)
        arc[start angle=\bA, end angle=0, radius=\rr]
        -- ($(B)+(0:\rr)+(0,-0.11)$);}
  \dt{O}
  \node[vlab] at ($(O)+(110:\R+0.30)$) {$A$};
  \node[vlab] at ($(O)+(180:\R+0.30)$) {$B$};
  \node[vlab] at ($(O)+(36:\R+0.30)$) {$C$};
  \node[vlab] at ($(O)+(0:\R+0.30)$) {$Q$};
  % The book sets O above-left of the centre dot, but the wedge there -- between
  % the dashed diameter and chord BC -- is only 0.42 r deep at the centre and
  % will not hold a footnotesize letter clear of both.  Below-left is the one
  % free quadrant; the dot itself now fixes which point the letter names.
  \node[vlab] at ($(O)+(250:0.32)$) {$O$};
\end{tikzpicture}
\figcap{12--22}
\end{minipage}
\end{bkfigure}}

%---------------------------------------------------------------- Fig 12-23
% Angle inscribed in a semicircle: BC is a diameter, so angle BAC is right.
\newcommand{\FIGXIITWENTYTHREE}{\begin{bkfigure}[\linewidth]
\begin{tikzpicture}[scale=1.300]
  \coordinate (O) at (0,0);
  \def\R{1.25}
  \coordinate (B) at ($(O)+(180:\R)$);
  \coordinate (C) at ($(O)+(0:\R)$);
  \coordinate (A) at ($(O)+(105:\R)$);
  \draw[given] (O) circle[radius=\R];
  \draw[given] (B) -- (C);
  \draw[given] (B) -- (A);
  \draw[given] (A) -- (C);
  \node[vlab] at ($(O)+(105:\R+0.30)$) {$A$};
  \node[vlab] at ($(O)+(180:\R+0.30)$) {$B$};
  \node[vlab] at ($(O)+(0:\R+0.30)$) {$C$};
  \node[vlab] at ($(O)+(265:0.34)$) {$O$};
\end{tikzpicture}
\figcap{12--23}
\end{bkfigure}}

%---------------------------------------------------------------- Fig 12-24
% Split out of the old \FIGXIITWENTYFOURTWENTYFIVE: 12-24 belongs to Exercise 7
% of Group 12-11 and 12-25 to Exercise *9, so they sit inline after different
% exercises.
% 12-24: AB and CD are parallel secants.  Their offsets from the centre are
%   DELIBERATELY unequal (0.55 r above, 0.68 r below -- both measured off the
%   plate).  Equal offsets would make AB congruent to CD, a coincidence the
%   exercise never grants; arc AC = arc BD holds for ANY pair of parallel
%   chords, which is exactly what Exercise 7 asks the student to prove.
% 12-25: AC is tangent at B (perpendicular to OB by construction); D and E are
%   placed symmetrically about OB, which makes DE parallel to AC and the arcs
%   BD and BE equal -- both facts the exercise's hint needs.
\newcommand{\FIGXIITWENTYFOUR}{\begin{bkfigure}[\linewidth]
\begin{tikzpicture}[scale=1.162]
  \coordinate (O) at (0,0);
  \def\R{1.25}
  \coordinate (A) at (-1.0440, 0.6875);
  \coordinate (B) at ( 1.0440, 0.6875);
  \coordinate (C) at (-0.9165,-0.8500);
  \coordinate (D) at ( 0.9165,-0.8500);
  \draw[given] (O) circle[radius=\R];
  \draw[given] (-1.52, 0.6875) -- (1.52, 0.6875);
  \draw[given] (-1.44,-0.8500) -- (1.44,-0.8500);
  \node[vlab] at ($(A)+(112:0.36)$) {$A$};
  \node[vlab] at ($(B)+(68:0.36)$) {$B$};
  \node[vlab] at ($(C)+(248:0.36)$) {$C$};
  \node[vlab] at ($(D)+(292:0.36)$) {$D$};
\end{tikzpicture}
\figcap{12--24}
\end{bkfigure}}

%---------------------------------------------------------------- Fig 12-25
% AC is tangent at B (perpendicular to OB by construction); D and E are placed
% symmetrically about OB, which makes DE parallel to AC and the arcs BD and BE
% equal -- both facts the exercise's hint needs.
\newcommand{\FIGXIITWENTYFIVE}{\begin{bkfigure}[\linewidth]
\begin{tikzpicture}[scale=1.221]
  \coordinate (O) at (0,0);
  \def\R{1.05}
  \coordinate (B) at ($(O)+(133:\R)$);
  \coordinate (A) at ($(B)+(223:1.22)$);
  \coordinate (C) at ($(B)+(43:1.06)$);
  \coordinate (D) at ($(O)+(199:\R)$);
  \coordinate (E) at ($(O)+(67:\R)$);
  \draw[given] (O) circle[radius=\R];
  \draw[given] (A) -- (C);
  \draw[aux] (B) -- (D);
  \draw[aux] (B) -- (E);
  \draw[aux] (D) -- (E);
  \node[vlab] at ($(A)+(223:0.34)$) {$A$};
  \node[vlab] at ($(O)+(133:\R+0.36)$) {$B$};
  \node[vlab] at ($(C)+(43:0.34)$) {$C$};
  \node[vlab] at ($(O)+(199:\R+0.34)$) {$D$};
  \node[vlab] at ($(O)+(67:\R+0.34)$) {$E$};
\end{tikzpicture}
\figcap{12--25}
\end{bkfigure}}

%---------------------------------------------------------------- Fig 12-26
% Split out of the old \FIGXIITWENTYSIXTWENTYSEVEN: 12-26 is first cited by
% Exercise *10 of Group 12-11 and 12-27 by Exercise *12, so each stands inline
% after its own item.
% 12-26: two secants from the external point A.  A is DERIVED as the crossing
%   of lines CB and DE, so A, B, C and A, E, D are exactly collinear.  The
%   auxiliary chord is B--D, NOT C--D: the book draws BD, and both exercises
%   that lean on this figure need triangle ABD -- Ex *10 gets angle A as the
%   exterior angle of ABD at B, and Ex *14 proves ABD ~ AEC.  The upper secant
%   runs on past C, as in the book, to show it is a line and not a segment.
\newcommand{\FIGXIITWENTYSIX}{\begin{bkfigure}[\linewidth]
\begin{tikzpicture}[scale=0.999]
  \coordinate (O) at (0,0);
  \def\R{1.1}
  \coordinate (B) at ($(O)+(151.6:\R)$);
  \coordinate (C) at ($(O)+(79.5:\R)$);
  \coordinate (D) at ($(O)+(-31:\R)$);
  \coordinate (E) at ($(O)+(213:\R)$);
  \coordinate (A) at (intersection of C--B and D--E);
  \draw[given] (O) circle[radius=\R];
  \draw[given] (A) -- ($(A)!1.12!(C)$);
  \draw[given] (A) -- (D);
  \draw[given] (B) -- (D);
  \node[vlab] at ($(A)+(186:0.34)$) {$A$};
  \node[vlab] at ($(B)+(88:0.50)$) {$B$};
  \node[vlab] at ($(O)+(95:\R+0.36)$) {$C$};
  \node[vlab] at ($(E)+(265:0.44)$) {$E$};
  \node[vlab] at ($(O)+(-31:\R+0.34)$) {$D$};
\end{tikzpicture}
\figcap{12--26}
\end{bkfigure}}

%---------------------------------------------------------------- Fig 12-27
% Chords AC and BD crossing at E, again derived as an intersection; BC dashed
% is the hint line for Exercise *12.
\newcommand{\FIGXIITWENTYSEVEN}{\begin{bkfigure}[\linewidth]
\begin{tikzpicture}[scale=1.140]
  \coordinate (O) at (0,0);
  \def\R{1.1}
  \coordinate (A) at ($(O)+(190:\R)$);
  \coordinate (B) at ($(O)+(135:\R)$);
  \coordinate (C) at ($(O)+(10:\R)$);
  \coordinate (D) at ($(O)+(-68:\R)$);
  \coordinate (E) at (intersection of A--C and B--D);
  \draw[given] (O) circle[radius=\R];
  \draw[given] (A) -- (C);
  \draw[given] (B) -- (D);
  \draw[aux] (B) -- (C);
  \node[vlab] at ($(O)+(190:\R+0.34)$) {$A$};
  \node[vlab] at ($(O)+(155:\R+0.36)$) {$B$};
  \node[vlab] at ($(O)+(10:\R+0.34)$) {$C$};
  \node[vlab] at ($(O)+(-68:\R+0.34)$) {$D$};
  \node[vlab] at ($(E)+(247:0.36)$) {$E$};
\end{tikzpicture}
\figcap{12--27}
\end{bkfigure}}

%---------------------------------------------------------------- Fig 12-28
% Two chords crossing at O.  2026-08-17: this used to be SOLVED so that the
% drawing reproduced the numbers quoted in Exercise 15 (AO:OB = 12:4 = 3 and
% CO:OD = 8:6 = 4/3) -- which handed the student the answer off a ruler.
% Caleb asked for the book's own schematic cut instead.  Measured off the
% plate, book p.224: A 198, B 31, C 105, D -18, giving AO:OB = 2.76 and
% CO:OD = 1.11.  Neither matches the quoted ratio, which is the point; what
% the figure still asserts exactly is the STRUCTURE -- four points on the
% circle, two chords, one crossing strictly inside.
\newcommand{\FIGXIITWENTYEIGHT}{\begin{bkfigure}[\linewidth]
\begin{tikzpicture}[scale=1.495]
  \coordinate (Ctr) at (0,0);
  \def\R{1.1}
  \coordinate (A) at ($(Ctr)+(198:\R)$);
  \coordinate (B) at ($(Ctr)+(31:\R)$);
  \coordinate (C) at ($(Ctr)+(105:\R)$);
  \coordinate (D) at ($(Ctr)+(-18:\R)$);
  \coordinate (X) at (intersection of A--B and C--D);
  \draw[given] (Ctr) circle[radius=\R];
  \draw[given] (A) -- (B);
  \draw[given] (C) -- (D);
  \node[vlab] at ($(Ctr)+(198:\R+0.30)$) {$A$};
  \node[vlab] at ($(Ctr)+(31:\R+0.30)$) {$B$};
  \node[vlab] at ($(Ctr)+(105:\R+0.30)$) {$C$};
  \node[vlab] at ($(Ctr)+(-18:\R+0.30)$) {$D$};
  \node[vlab] at ($(X)+(250:0.36)$) {$O$};
\end{tikzpicture}
\figcap{12--28}
\end{bkfigure}}

%---------------------------------------------------------------- Fig 12-29
% Two secants from an external point A.  2026-08-17: this was SOLVED from the
% numbers in Exercise *16 (AD = 6, DE = 10, AB = 8, hence BC = 4), which is
% precisely what Caleb asked us to stop doing -- a scale drawing answers the
% exercise for the student.  It is now the book's own schematic cut, measured
% off the plate on p.225: D 186, E 47, B 203, C -21, and A DERIVED as the
% crossing of the two secants, so collinearity is exact rather than plotted.
% Two further faults from the same photo pair, both fixed by these angles:
%   * the lower secant used to graze the circle -- B and C landed close
%     together on the bottom arc.  The book runs BC as a long, near-horizontal
%     chord across almost the whole circle, with C on the RIGHT.
%   * A sat about 1.5 diameters out with a 25-30 degree wedge.  The book keeps
%     A close in (about 0.6 r clear) and the wedge compact.
% What the drawing still asserts exactly: four points on the circle, both
% secants straight through A, D between A and E, B between A and C.
\newcommand{\FIGXIITWENTYNINE}{\begin{bkfigure}[\linewidth]
\begin{tikzpicture}[scale=1.140]
  \coordinate (O) at (0,0);
  \def\R{1.4}
  \coordinate (D) at ($(O)+(186:\R)$);
  \coordinate (E) at ($(O)+(47:\R)$);
  \coordinate (B) at ($(O)+(203:\R)$);
  \coordinate (C) at ($(O)+(-21:\R)$);
  \coordinate (A) at (intersection of D--E and B--C);
  \draw[given] (O) circle[radius=\R];
  \draw[given] (A) -- (E);
  \draw[given] (A) -- (C);
  \node[vlab] at ($(A)+(196:0.34)$) {$A$};
  \node[vlab] at ($(D)+(120:0.32)$) {$D$};
  \node[vlab] at ($(B)+(255:0.32)$) {$B$};
  \node[vlab] at ($(O)+(47:\R+0.32)$) {$E$};
  \node[vlab] at ($(O)+(-21:\R+0.32)$) {$C$};
\end{tikzpicture}
\figcap{12--29}
\end{bkfigure}}

%---------------------------------------------------------------- Fig 12-30
% Tangent plus secant from an external point.  2026-08-17: likewise reverted
% from the solved cut (AC = 8, CD = 10, AB = 12) to the book's schematic.  The
% two structural facts the exercises actually lean on are now built in, and
% both were broken before:
%   * AB is TANGENT at B.  The old cut met the circle at a visibly oblique
%     angle to the radius, so it read as a clipped secant -- and Ex *17's whole
%     argument is the tangent case.  B is now derived from the tangent-length
%     construction, so OB is perpendicular to AB exactly.
%   * ACD is a true DIAMETER, level with A, as the book draws it; the old cut
%     sloped it down the right-hand side and it missed the centre entirely.
% With A on the line of the diameter, AB^2 = AC * AD then holds identically --
% but the drawn ratio AC:CD is 0.43, not the exercise's 0.8, so nothing is
% given away.
\newcommand{\FIGXIITHIRTY}{\begin{bkfigure}[\linewidth]
\begin{tikzpicture}[scale=1.140]
  \coordinate (O) at (0,0);
  \def\R{1.4}
  \def\OA{2.59}
  \coordinate (A) at (-\OA, 0);
  \coordinate (C) at ($(O)+(180:\R)$);
  \coordinate (D) at ($(O)+(0:\R)$);
  % tangency: the radius to B makes angle arccos(r/|OA|) with OA, which is the
  % construction, not a guess -- so OB is exactly perpendicular to AB.
  \pgfmathsetmacro{\Bang}{180-acos(\R/\OA)}
  \coordinate (B) at ($(O)+(\Bang:\R)$);
  \draw[given] (O) circle[radius=\R];
  \draw[given] (A) -- (B);
  \draw[given] (A) -- (D);
  \draw[aux] (B) -- (C);
  \draw[aux] (B) -- (D);
  \node[vlab] at ($(A)+(186:0.32)$) {$A$};
  \node[vlab] at ($(B)+(100:0.32)$) {$B$};
  \node[vlab] at ($(C)+(310:0.34)$) {$C$};
  \node[vlab] at ($(O)+(0:\R+0.32)$) {$D$};
\end{tikzpicture}
\figcap{12--30}
\end{bkfigure}}

%---------------------------------------------------------------- Fig 12-31
% A central right angle; the dashed chord is the segment of length sqrt(2) r
% against which the arc is compared.
\newcommand{\FIGXIITHIRTYONE}{\begin{bkfigure}[\linewidth]
\begin{tikzpicture}[scale=1.300]
  \coordinate (O) at (0,0);
  \def\R{1.25}
  \coordinate (T) at ($(O)+(90:\R)$);
  \coordinate (S) at ($(O)+(0:\R)$);
  \draw[given] (O) circle[radius=\R];
  \draw[given] (O) -- (T);
  \draw[given] (O) -- (S);
  \draw[aux] (T) -- (S);
  \node[slab] at ($(O)+(90:0.62)+(-0.26,0)$) {$r$};
\end{tikzpicture}
\figcap{12--31}
\end{bkfigure}}


% ---- local macros (hoist into book preamble) ----
% Arc accent, as the book sets it: a frown over the endpoint letters.
\newcommand{\XIIarc}[1]{\overset{\raisebox{-0.25ex}{\scalebox{1.15}[0.85]{$\frown$}}}{#1}}
% Equation numbers as (12--1), matching the book's en-dash style.
\renewcommand{\theequation}{\thechapter--\arabic{equation}}
% Numbered remarks. The book sets these as italic "Remark" + roman numeral,
% with NO period after the word (e.g. "Remark 1."), unlike the unnumbered
% \begin{remark} in brumfiel.sty, which prints "Remark." Verified against
% book pp. 218, 221.
\newenvironment{XIIremarkn}[1]{\par\medskip\noindent\emph{Remark} #1.\ }%
  {\par\medskip}
% Footnotes marked with * rather than a digit (the book's convention).
\newcommand{\XIIsymfoot}[1]{%
  \begingroup
  \renewcommand{\thefootnote}{\fnsymbol{footnote}}%
  \setcounter{footnote}{0}%
  \footnote{#1}%
  \endgroup}
% ---- end local macros ----

\begin{document}
\setcounter{chapter}{12}

\chapter*{\raggedright\Huge MEASUREMENT OF\\ANGLES AND ARCS}
\addcontentsline{toc}{chapter}{12.\ Measurement of Angles and Arcs}
\markboth{MEASUREMENT OF ANGLES AND ARCS}{}
\vspace{-1em}\noindent\rule{\linewidth}{0.6pt}\vspace{1em}
\figurekey

%========================================================= 12-1
\section*{12--1\quad INTRODUCTION}
\addcontentsline{toc}{section}{12--1\quad Introduction}

In the first chapter of this book you used a protractor to measure
representations of angles. Finally, at this late stage, we prove that the
process of angle measurement can be made precise in our geometry. In order to
discuss angle measurement, it is convenient first to define the measure of an
arc of a circle. This definition is quite like the definition of circumference
and really only makes precise the intuitive idea of \emph{length of arc} that
you already have. We needed limits to discuss circumference, and of course we
also need limits in order to discuss arc length.

Once we know what is meant by arc length, it will be relatively easy to discuss
measurement of angles. This discussion will go quite like Chapter 5 on
measurement of length of segments.

One reason that a precise description of measurement of angles in degrees is
difficult is that describing a unit angle is troublesome. (Is there an angle of
one degree?) Furthermore, how do we know that it is possible to divide a degree
into 60 congruent angles (minutes) and a minute into successively smaller
pieces? This chapter shows how to resolve these difficulties. Incidentally, it
should become quite clear that what we choose as the \emph{unit angle} is quite
arbitrary, and it is important for you to understand how measurements differ
when different units are used. The idea that ties all this together is that of
\emph{length of circular arc}.

%========================================================= 12-2
\section*{12--2\quad CIRCULAR ARCS}
\addcontentsline{toc}{section}{12--2\quad Circular arcs}

We have already described the interior and exterior of an angle (Chapter 3),
and we use Definition 3--10 to describe an arc of a circle.

\begin{definition}
If $\ang{AOB}$ is a central angle of a circle, \emph{the arc $\XIIarc{AB}$ of
the circle} subtended by $\ang{AOB}$ is defined to be the set of points of the
circle which are \emph{interior} to $\ang{AOB}$ or are on the angle. The set of
points of the circle which are \emph{exterior} to or on the angle is also called
\emph{an arc of the circle} (for example, the arc $\XIIarc{ABC}$ in
Fig.~12--1). If the central angle is a straight angle (like $\ang{COB}$ in
Fig.~12--1), then it has no interior, but the set of points of the circle which
are on one side of $CB$ or on $CB$ is called a \emph{semicircular arc} and is
said to be subtended by $\ang{COB}$. Finally, when we speak simply of an
\emph{arc of a circle}, we mean an arc subtended by a central angle or the arc
exterior to some central angle.
\end{definition}

\FIGXIIONE

\exercisegroup{12--1}
\begin{multicols}{2}
\begin{exlist}
\item In Fig.~12--2, show that if $\ang{AOB} > \ang{COB}$, with $C$ in the
      interior of $\ang{AOB}$, then every point of the arc $\XIIarc{CB}$ is a
      point of the arc $\XIIarc{AB}$.
\item If the arc $\XIIarc{AB}$ is subtended by the central angle $AOB$, define
      the \emph{midpoint of the arc}, and show that such a point exists.
\item Use the result of Exercise 2 to prove that there are infinitely many
      points on the arc.
\item In Fig.~12--3, the arc $\XIIarc{EF}$ is subtended by $\ang{EOF}$. Consider
      the segment $EF$ and any point $P$ on $EF$. Show that there is, on the ray
      $OP$, a point $Q$, which is on the arc $\XIIarc{EF}$. Show that if
      $P \neq P'$, then $Q \neq Q'$. How many points are there on the arc
      $\XIIarc{EF}$?
\item Referring to Fig.~12--4, suppose that an arc is exterior to a central
      angle $AOB$. Show how to define the midpoint of the arc. Why is
      $\ang{AOC} \cong \ang{BOC}$?
\end{exlist}
\end{multicols}

\FIGXIITWOTHREE

\FIGXIIFOUR

%========================================================= 12-3
\section*{12--3\quad ARC LENGTH}
\addcontentsline{toc}{section}{12--3\quad Arc length}

In defining the length of a segment, we first had to choose a unit of length.
So also in defining arc length, it will be understood that we have already
chosen a segment of unit length. We shall also suppose that we have defined the
midpoint of an arc as in Exercise Group 12--1.

We describe how to construct a regular inscribed polygonal path in an arc, and
indeed, polygonal paths of 2, 4, 8, 16,\ldots\ sides.

In Fig.~12--5, let $AB$ be a circular arc, and let $C$ be the midpoint of the
arc. (See Exercise Group 12--1.) The segments $AC$ and $CB$ form a regular
inscribed polygonal path from $A$ to $B$, with two sides. There are now formed
two congruent angles, $\ang{AOC}$ and $\ang{COB}$. Each of these angles subtends
an arc. Let $D$ and $E$ be the midpoints of these arcs. We can now consider the
regular inscribed polygonal path of four sides, $ADCEB$. We then consider each
of the arcs $\XIIarc{AD}$, $\XIIarc{DC}$, $\XIIarc{CE}$, and $\XIIarc{EB}$ and
their midpoints $F$, $G$, $H$, and $I$. Thus we obtain a regular inscribed path
$AFDGCHEIB$ from $A$ to $B$, having eight sides. Continuing in this manner, for
each positive integer we obtain a regular inscribed polygonal path from $A$ to
$B$. That is to say, first we have a path of two sides, next a path of four
sides, then a path of eight sides, and so on.

\FIGXIIFIVESIX

The following facts are clear from Fig.~12--6, and a complete proof is not
difficult to give, although several steps are needed. We shall use these facts
without proof. The \emph{length} of the circumscribed path in Fig.~12--6 will be
denoted by $L$. The length of a regular inscribed path of 2, 4, 8,
16,\ldots, sides will be denoted by $l_n$, where $n$ is used to denote the
number of segments in the path.

(a) \emph{For every $n$, $l_n < L$}; that is to say, the length of each regular
inscribed path is less than the length of any circumscribed path.

(b) The length of each regular inscribed path of 2, 4, 8, 16,\ldots\ sides is
less than the one that follows, that is, $l_n < l_{2n}$.

(c) From (a) and (b) and properties of real numbers, it follows that
$\lim \{l_n\}$ exists.

We are now able to define length of arc.

\begin{definition}
Given a circular arc, let $l_n$ be the length of a regular inscribed polygonal
path of 2, 4, 8, 16, 32,\ldots, sides inscribed in the arc. The \emph{length $l$
of the arc} is defined to be the following limit: $l = \lim \{l_n\}$.
\end{definition}

\exercisegroup{12--2}
\begin{multicols}{2}
\begin{exlist}
\item A regular hexagon is inscribed in a circle of radius 1, as in
      Fig.~12--7. Show that the length of the arc $\XIIarc{ACB}$ is greater
      than 2.
\item In a circle of radius 2, what is the length of arc subtended by a right
      angle? by a straight angle?
\item In Fig.~12--8, the central angles $AOB$ and $BOC$ are congruent, and
      $\ang{AOC}$ is a right angle. What is the length of the arc $\XIIarc{AB}$
      if $\sg{AO} = 5$?
\item[\stex 4.] Prove that the length of any arc is a positive number which is
      greater than the length of the line segment determined by the endpoints of
      the arc.
\end{exlist}
\end{multicols}

\FIGXIISEVENEIGHT

\begin{theorem}
If, in a circle, two central angles are congruent, they subtend arcs of equal
length.
\end{theorem}

\noindent\emph{Proof.} Because the angles are congruent, the inscribed polygons
at any stage are also congruent. Thus the lengths of the corresponding
inscribed polygonal paths are the same. Hence the arc lengths are the same.

\begin{remark}
In the same way it follows that the arcs exterior to the central angles also
have the same length.
\end{remark}

Having defined arc length, we now would like to prove some more theorems about
arc lengths. Indeed, it would be possible to proceed at once to this task, but
if we were to do so, our proofs would be quite troublesome. In constructing
later proofs it will be most convenient for us to use a theorem (12--2 below)
which tells us that in approximating to the circle by inscribed polygons, we
need \emph{not} restrict ourselves to regular polygons. The proof of this
``helping'' theorem is rather hard to make and involves greater knowledge of
limits than is presupposed here. Nevertheless, the geometric significance of the
theorem should be clear, and it should appeal to your intuition. We shall use
this theorem without proof, so that actually we treat it like a postulate.

\begin{theorem}
Let $\{P_n\}$ be any sequence of polygonal paths inscribed in the arc
$\XIIarc{AB}$ of a circle (Fig.~12--9), and let $\{l_n\}$ be the sequence of
lengths of these polygonal paths; let $s_n$ be the length of the longest side of
the path $P_n$. If $\lim \{s_n\} = 0$, then $\lim \{l_n\} = $ length
$\XIIarc{AB} = l$.
\end{theorem}

\begin{remark}
Evidently, when the length of the longest side of $P_n$ is small, the number of
sides, $n$, must be large. We can now use this theorem to prove some facts about
arc length.
\end{remark}

\begin{theorem}
If $\XIIarc{AB}$ and $\XIIarc{BC}$ (Fig.~12--10) are adjacent arcs (that is, if
they have only point $B$ in common), then length $\XIIarc{ABC} = $ length
$\XIIarc{AB} + $ length $\XIIarc{BC}$.
\end{theorem}

\FIGXIININETEN

\noindent\emph{Proof.} The length of $\XIIarc{ABC}$ is, according to
Theorem 12--2, the limit of lengths of inscribed polygonal paths from $A$ to
$C$. Let us always choose these paths so that $B$ is one of the vertices of the
path. The path from $A$ to $C$ can then be separated into a path from $A$ to $B$
and a path from $B$ to $C$. Let $l'_n$ be the length of the path from $A$ to
$B$, and $l''_n$ be the length of the path from $B$ to $C$. Then, if $l_n$ is
the length of the path from $A$ to $C$, we have $l_n = l'_n + l''_n$.

Now we consider sequences of paths $\{P'_n\}$ and $\{P''_n\}$, with the
sequences of lengths of longest sides approaching zero. Thus,
\begin{equation}
\left.
\begin{aligned}
\lim \{l_n\} &= \text{length } \XIIarc{ABC},\\
\lim \{l'_n\} &= \text{length } \XIIarc{AB},\\
\lim \{l''_n\} &= \text{length } \XIIarc{BC}.
\end{aligned}
\right\}
\end{equation}
But from a theorem on limits, we have
\begin{equation}
\lim \{l'_n + l''_n\} = \lim \{l'_n\} + \lim \{l''_n\}.
\end{equation}
From Eqs.~(12--1) and (12--2), we see that
\begin{equation}
\text{length } \XIIarc{ABC} = \text{length } \XIIarc{AB} + \text{length }
\XIIarc{BC}.
\end{equation}
This completes the proof. The theorem might be stated briefly as ``the length of
the `sum' is equal to the sum of the lengths.''

\FIGXIIELEVENTWELVE

\begin{theorem}
If, in a circle, one central angle is greater than another central angle, the
greater angle subtends the longer arc.
\end{theorem}

\noindent\emph{Proof.} Since congruent angles subtend arcs of equal length, we
may suppose that the central angles are as in Fig.~12--11, and that
$\ang{AOB} > \ang{COB}$. The remainder of the proof (using Theorem 12--3) is
left to the student.

\exercisegroup{12--3}
\begin{multicols}{2}
\begin{exlist}
\item Apply Theorem 12--3 to the arcs subtended by a straight angle and two
      adjacent right angles.
\item State the form Theorem 12--3 would have for three adjacent arcs; for four
      adjacent arcs; for $n$ adjacent arcs.
\item Two adjacent arcs in a circle of radius $r$ have lengths $\pi r/3$ and
      $2\pi r/3$. What angle subtends the two arcs together?
\item Two arcs on a circle of radius $r$ have lengths $r$ and $\pi r/3$. Which
      is subtended by the greater central angle?
\end{exlist}
\end{multicols}

\refstepcounter{theorem}
\par\medskip\noindent\textbf{Theorem \thetheorem} Converse to Theorem 12--4.\
\emph{If one arc of a circle is longer than a second arc of the circle, the
longer arc is subtended by the greater central angle.}\par\medskip

The proof is left to the student.

\begin{theorem}
Two central angles of a circle are congruent if and only if they subtend arcs of
equal length.
\end{theorem}

The proof is left to the student.

%========================================================= 12-4
\section*{12--4\quad ARCS ON DIFFERENT CIRCLES}
\addcontentsline{toc}{section}{12--4\quad Arcs on different circles}

So far we have compared arcs on the same circle. We now compare the lengths of
arcs on different circles. Before studying Theorem 12--7, try to discover the
relation between the lengths of two arcs in different circles if they are
subtended by congruent angles.

\begin{theorem}
If $\XIIarc{AB}$ and $\XIIarc{A'B'}$ are arcs on two circles of radii $r$ and
$r'$, respectively, and are subtended by congruent angles, then the arc lengths
are proportional to the radii (Fig.~12--12). That is, if
$\ang{AOB} \cong \ang{A'O'B'}$, then
\[
\frac{\text{length } \XIIarc{AB}}{\text{length } \XIIarc{A'B'}} =
\frac{r}{r'}, \qquad \text{or} \qquad
\frac{\text{length } \XIIarc{AB}}{r} =
\frac{\text{length } \XIIarc{A'B'}}{r'}.
\]
\end{theorem}

\noindent\emph{Proof.} The proof depends upon three things: (1) the definition
of arc length, (2) a theorem on limits, and (3) properties of similar triangles.

Suppose that $APQB$ is a polygonal path inscribed in the arc $AB$. Then there is
a corresponding polygonal path $A'P'Q'B'$ inscribed in the arc $A'B'$ such that
corresponding central angles are congruent, that is,
$\ang{AOP} \cong \ang{A'O'P'}$, $\ang{POQ} \cong \ang{P'O'Q'}$, etc. (Why?)
Hence, there is a one-to-one correspondence between polygonal paths inscribed in
the two arcs. Let $L_n$ be the length of a polygonal path of $n$ sides inscribed
in $\XIIarc{AB}$, and $L'_n$ be the length of the corresponding polygonal path
inscribed in $\XIIarc{A'B'}$. Because the corresponding central angles in the
two polygons are congruent, for example $\ang{AOP} \cong \ang{A'O'P'}$, the
triangles formed are similar. (Why?) Since $\sg{AP}/r = \sg{A'P'}/r'$, etc., we
have $L_n/r = L'_n/r'$ and $L_n/L'_n = r/r'$.

From the definition of length of arc, length $\XIIarc{AB} = \lim \{L_n\}$ and
length $\XIIarc{A'B'} = \lim \{L'_n\}$. From a theorem on limits,
\[
\frac{\text{length } \XIIarc{AB}}{\text{length } \XIIarc{A'B'}} =
\frac{\lim \{L_n\}}{\lim \{L'_n\}} = \lim \left\{\frac{L_n}{L'_n}\right\} =
\frac{r}{r'}.
\]
This completes the proof.

\exercisegroup{12--4}
\begin{multicols}{2}
\begin{exlist}
\item A central angle $B$ subtends an arc of length 5 in.\ on a circle of radius
      4 in. A central angle $B' \cong \ang{B}$ subtends an arc of length 10 in.\
      on another circle. What is the radius of the second circle?
\item A central angle subtends an arc of 2 in.\ on a circle of radius 5 in. What
      would be the length of an arc subtended by this same angle on a circle
      with the same center but with a radius of 2 ft?
\item The radii of two circles have a ratio of 3 to 5. If a central angle
      subtends an arc of length 2.454 on the first circle, what length arc would
      a central angle that is congruent to the first angle subtend on the second
      circle?
\end{exlist}
\end{multicols}

%========================================================= 12-5
\section*{12--5\quad MORE ABOUT ARC LENGTH}
\addcontentsline{toc}{section}{12--5\quad More about arc length}

In the preceding sections we learned that every arc has associated with it a
number called its length. We shall not consider an exhaustive proof of the
converse theorem: \emph{For every positive number, less than the circumference,
there is an arc having this number as its length.} The complete proof of this
theorem requires a rather difficult argument using limits, so we will limit our
proof to certain special numbers. It will then be clear that the theorem is
true.

In proving the theorem, we confine ourselves to arcs of length not more than one
quarter of the circumference, that is to arcs which are subtended by acute
angles or right angles.

\begin{theorem}
Suppose that $L$ is a positive number less than or equal to one fourth of the
circumference of a circle of radius $r$ $(L \leqq \pi r/2)$.\XIIsymfoot{Since
the circumference of a circle is $2\pi r$, it follows that one fourth of the
circumference is $\frac{1}{4} \times 2\pi r$, or $\pi r/2$.} Then there is an
arc of the circle whose length is $L$.
\end{theorem}

\noindent``\emph{Proof}.'' If $L = \pi r/2$, then the required arc is subtended
by a right angle. If $L = \frac{1}{2} \cdot \pi r/2$, then bisecting a central
right angle will give two arcs $\XIIarc{AC}$ and $\XIIarc{BC}$, as in
Fig.~12--13, of length $L$. If $L = \frac{3}{4} \cdot \pi r/2$, then the arc
$AD$ of Fig.~12--13 has length $L$, where $OD$ bisects $\ang{BOC}$. Similarly,
if $OE$ bisects $\ang{DOC}$, then length $\XIIarc{AE} = \frac{5}{8} \cdot
\pi r/2$.

Thus we see that if $L$ is a certain special fractional part of $\pi r/2$, then
there is an arc of length $L$. (Try to describe these fractions.) The fractions
for which arcs are easy to construct are 1/2, 1/4, 3/4, 1/8, 3/8, 5/8, 7/8,
1/16,\ldots. In other words, the arc is easy to construct if its length is given
simply in terms of halves, quarters, eighths, etc., of a quarter of the
circumference.

\FIGXIITHIRTEENFOURTEEN

You are familiar with measuring lengths with a ruler whose scale is divided into
halves, quarters, eighths, sixteenths, etc. Hence, if the number $L$ is not
exactly expressible in halves, quarters, eighths, etc., we can still construct
an arc whose length is very close to $L$.

In Fig.~12--14, $L$ appears to be a little greater than $\frac{5}{8} \cdot
\pi r/2$. It is plausible therefore, that for every $L < \pi r/2$, there is
exactly one arc $AP$, as in the figure, with length $\XIIarc{AP} = L$.

\exercisegroup{12--5}
\begin{multicols}{2}
\begin{exlist}
\item In this section we gave a list of fractions whose associated arcs are easy
      to construct. The list is computed according to a definite rule. Continue
      the list to 20 more terms.
\item On a circle of radius $r$, construct arcs of length
      $7/8 \cdot \pi r/2$; $9/16 \cdot \pi r/2$; $7/4 \cdot \pi r/2$.
\item Suppose that there is an arc of length 3 on a circle of radius 2 in. What
      part (as a decimal) of $\pi r/2$ is 3? Show how to approximate this
      decimal by a sum of halves, quarters, eighths, etc. Use the approximation
      3.14 for $\pi$, and then rework using 3.1416.
\item An arc of length $r$ on a circle of radius $r$ is what part of the
      circumference? of one quarter of the circumference?
\item On a circle of radius 32 in.\ is an arc of length 22 in. What part of a
      right angle does it subtend? Use $\pi = 22/7$. Give a construction
      approximating this arc.
\end{exlist}
\end{multicols}

%========================================================= 12-6
\section*{12--6\quad MEASUREMENT OF ANGLES}
\addcontentsline{toc}{section}{12--6\quad Measurement of angles}

In the previous section we have in effect been measuring angles by measuring
arcs. In fact we have expressed the arc length as a fraction of $\pi r/2$ that
is as a fraction of one quarter of the circumference. In doing this we have
measured the corresponding central angle in terms of the right angle as a unit
of angle. In this section we shall define the measure of an angle in terms of
other units.

In deciding how to measure angles, we must first define the unit. This involves
choosing some angle whose measure is 1. The choice is quite arbitrary, but there
are some restrictions based on convenience. If the unit angle is small, the
measures of common angles will be large numbers; whereas if the unit angle is
large, the measures of common angles will be small numbers. It turns out that
the familiar degree is a convenient size unit for ordinary measurement, because
the numbers involved are a ``good'' size. Obviously, other unit angles could be
selected. A natural one would be obtained by dividing the circle into 100 arcs
of equal length, or perhaps the right angle into 100 congruent angles.

How does it happen that we use that rather strange unit, the degree? We measure
angles in degrees because the ancient Sumerians\XIIsymfoot{The Sumerians had a
highly developed civilization in Mesopotamia, in the valleys of the Tigris and
Euphrates rivers, more than 5000 years ago. Roughly contemporary with ancient
Egypt, the Sumerian culture discovered much about arithmetic, geometry, and
astronomy. We inherit many things from these gifted people. Their method of
counting was based on 60 rather than our base of 10. Thus we divide the hour
into 60 minutes, and the minute into 60 seconds. They also invented place value
in number notation, although they had no ``zero.'' Many of our words go back to
Sumerian roots. For more Sumerian history, see \emph{A History of Science}, by
George Sarton, Harvard University Press, 1952, or your encyclopedia.} (2000
B.C.\ and earlier) did. They divided the circle into 360 ``equal parts,''
perhaps because they thought at one time that the year was 360 days long!

\begin{definition}
On a circle of radius $r$, an arc of length $1/360 \times 2\pi r$ is subtended
by a central angle of one degree $(1^\circ)$.
\end{definition}

\begin{XIIremarkn}{1}
The definition of degree is independent of the radius $r$, by
Theorem 12--7.
\end{XIIremarkn}

\begin{XIIremarkn}{2}
Any angle congruent to an angle of $1^\circ$ is also an angle of
$1^\circ$, by Theorem 12--6.
\end{XIIremarkn}

\exercisegroup{12--6}
\begin{multicols}{2}
\begin{exlist}
\item Write out the proof of Remark 1.
\item Write out the proof of Remark 2.
\item Define ``minute'' and ``second.''
\item Define a new unit of angle, which we will call the \emph{centile}, by
      dividing one quarter of the circumference into 100 equal arcs. Using this
      unit, what is the measure of a $30^\circ$ angle? a $45^\circ$ angle? a
      straight angle?
\item Divide the circumference of a circle into two arcs of equal length. What
      angle is subtended by each arc? Each arc is what part of the
      circumference? What is the measure in degrees of the central angle that is
      subtended by each arc?
\item Divide the circumference of a circle into four arcs of equal length. What
      kind of an angle is subtended by each of these arcs? Each arc is what part
      of the circumference? What is the measure in degrees of each central
      angle?
\item Consider an arc whose length is one sixth of the circumference. What is
      the measure in degrees of the central angle that is subtended by this arc?
\item Consider an arc whose length is one sixth of the length of a semicircular
      arc. What is the measure in degrees of the central angle that is subtended
      by this arc?
\item If a semicircular arc has length 8 in., and a second arc has length 2 in.,
      what part is the length of the second arc of the first? What is the
      measure in degrees of the central angle that is subtended by the second
      arc?
\item On a circle of radius 4, what is the length of an arc subtended by an
      angle of $1^\circ$? of $2^\circ$? of $30^\circ$?
\end{exlist}
\end{multicols}

\begin{definition}
If, on a circle of radius $r$, a central angle subtends an arc of length $L$,
then the measure of the angle in degrees is $L/\pi r \times 180$.
\end{definition}

\begin{XIIremarkn}{1}
This definition agrees with Definition 12--3 for angles of $1^\circ$.
\end{XIIremarkn}

\begin{XIIremarkn}{2}
By Theorem 12--7, the measure of an angle in degrees is independent of the
radius.
\end{XIIremarkn}

\begin{XIIremarkn}{3}
If $M$ is the measure of the central angle in degrees, we have from
Definition 12--4 that $L = M\pi r/180$.
\end{XIIremarkn}

\exercisegroup{12--7}
\begin{multicols}{2}
\begin{exlist}
\item If $\ang{AOB}$ and $\ang{BOC}$ are adjacent angles of $1^\circ$ each, show
      that the measure of $\ang{AOC}$ is $2^\circ$.
\item Show that a straight angle has measure $180^\circ$.
\item On a circle of radius 6, find the number of degrees in a central angle
      that subtends an arc of length 3; 6; $\pi$; $2\pi/3$.
\item On a circle of radius 9, find the length of an arc subtended by an angle
      of $30^\circ$; $180^\circ$; $60^\circ$; $(180/\pi)^\circ$.
\item Show that bisecting a right angle gives an angle whose measure is
      $45^\circ$.
\item Prove Remark 3 under Definition 12--4.
\end{exlist}
\end{multicols}

%========================================================= 12-7
\section*{12--7\quad COUNTING DEGREES AND\\MEASURING ANGLES}
\addcontentsline{toc}{section}{12--7\quad Counting degrees and measuring angles}

Measurement of angles is quite similar to the counting process which led to the
idea of length of a segment. In this section we give a description of this
process and find that it gives the same measure of an angle as is given by
Definition 12--4. Detailed proofs will be omitted. The following theorem shows
how counting applies.

\begin{theorem}
Suppose that $\ang{AOB}$ is an acute angle, and we consider successive arcs
subtending central angles of $1^\circ$ on the circle from $B$ to $A$ so that the
point $A$ is a point of the $(n + 1)$st arc (Fig.~12--15). Then the measure of
$\ang{AOB}$ in degrees is either equal to $n$, or $(n + 1)$, or is between $n$
and $(n + 1)$; $n \leqq$ measure of $\ang{AOB} \leqq n + 1$.
\end{theorem}

\FIGXIIFIFTEEN

\noindent\emph{Proof.} The length of the arc of $n^\circ$ is
$n \cdot \pi r/180$; the length of the arc of $(n + 1)^\circ$ is
$[(n + 1)\pi r]/180$. The length $L$ of the arc $\XIIarc{AB}$ is then between
these two numbers, by Theorem 12--5.
\begin{equation}
n \cdot \frac{\pi r}{180} \leqq L \leqq \frac{(n + 1)\pi r}{180}.
\end{equation}
By Definition 12--4, the measure of $\ang{AOB}$ in degrees is
$L \cdot 180/\pi r$. From Eq.~(12--4), we see that
\[
n \leqq L \cdot \frac{180}{\pi r} \leqq (n + 1),
\]
and this is what we wished to prove.

\begin{remark}
The measure to the nearest minute could be obtained by counting in a similar
way.
\end{remark}

Consider an angle $A'O'B'$ (Fig.~12--16) that is congruent to $\ang{AOB}$ of
Theorem 12--9. Consider also successive arcs subtending central angles of
$1^\circ$ on arc $B'A'$, as described for arc $BA$. If the point $A$ is on the
$(n + 1)$st arc, then the point $A'$ is on the $(n + 1)$st arc. This can be
shown using the addition and subtraction of angles postulate and the definition
of ``less than'' and ``greater than'' for angles. This proves the following
theorem:

\begin{theorem}
If two angles are congruent, they have equal measures.
\end{theorem}

The converse of Theorem 12--10 can be proved, using the fact that
$\ang{O} > \ang{O'}$, or $\ang{O} \cong \ang{O'}$, or $\ang{O} < \ang{O'}$.

\FIGXIISIXTEEN

\exercisegroup{12--8}
\begin{center}Give your answers in terms of $\pi$.\end{center}
\nopagebreak
\begin{multicols}{2}
\begin{exlist}
\item If the radius of a circle is 14, what is the length of the arc subtended
      by an angle of $30^\circ$?
\item On a circle of radius 5 in.\ is an arc of length 6 in. What is the measure
      in degrees of the central angle subtending the arc? What if the radius is
      20 in.? 6 in.?
\item A central angle of a circle subtends an arc of length equal to the radius.
      How many degrees in the angle?
\item An arc of length 8 in.\ on a circle of radius 6 in.\ subtends a central
      angle of how many degrees? an arc of length 4 in.? an arc of length
      6 in.?
\item On a circle of radius $\pi$, there is an arc of length $\pi^2/4$. How many
      degrees in the central angle?
\item An arc of length $\pi r/2$ is subtended on a circle of radius $r$ by what
      central angle?
\item Find the length of arc on a circle of radius $r$ subtended by an angle of
      $1^\circ$.
\end{exlist}
\end{multicols}

We have been measuring arcs by their lengths. It is convenient to measure them
by the measure in degrees of the central angles they subtend.

\begin{definition}
The measure in degrees of an arc subtended by a central angle is equal to the
measure of the central angle. The measure in degrees of an arc exterior to a
central angle is equal to 360 minus the measure of the central angle. (See
Fig.~12--17.)
\end{definition}

\FIGXIISEVENTEENEIGHTEEN

\exercisegroup{12--9}
\begin{multicols}{2}
\begin{exlist}
\item An arc of length $\pi$ in.\ on a circle of radius 3 in.\ is of how many
      degrees?
\item An arc of length 4 in.\ on a circle of radius 6 in.\ has how many degrees?
\item On a circle of radius 2 in., how long is an arc of $120^\circ$?
      $240^\circ$?
\item How many degrees has an arc of length 15 in.\ on a circle of radius 3?
\end{exlist}
\end{multicols}

%========================================================= 12-8
\section*{12--8\quad RADIAN MEASURE}
\addcontentsline{toc}{section}{12--8\quad Radian measure}

Besides the familiar unit of degree for measuring angles, there is another one
which is of great importance in advanced mathematics. The reasons for using it
cannot be given here, but the definition is quite simple, and you will meet it
again if you study trigonometry. The unit is called the \emph{radian}.

\begin{definition}
An angle of one radian is one which subtends an arc of length $r$ on a circle of
radius $r$ (Fig.~12--18).
\end{definition}

\begin{definition}
The measure of an angle in radians is equal to $L/r$, where $L$ is the length of
the arc subtended on a circle of radius $r$.
\end{definition}

\begin{XIIremarkn}{1}
Definition 12--7 agrees with Definition 12--6 for angles of 1 rad.
\end{XIIremarkn}

\begin{XIIremarkn}{2}
The measure in radians does not depend on the radius, by Theorem 12--7.
\end{XIIremarkn}

\exercisegroup{12--10}
\begin{multicols}{2}
\begin{exlist}
\item On a circle of radius 2 is an arc of length 4. How many radians are there
      in the central angle?
\item How many radians are there in a right angle? in a straight angle?
\item On a circle of radius 2, find the length of arc subtended by a central
      angle of $30^\circ$. How many radians are there in this angle?
\item How many degrees are there in one radian?
\item How many radians are there in one degree?
\item From Exercises 4 and 5, derive the formulas for changing measure in
      degrees to measure in radians, and vice versa.
\item What are the measures in radians of the following angles? $30^\circ$;
      $45^\circ$; $60^\circ$; $90^\circ$; $120^\circ$; $135^\circ$;
      $150^\circ$; $180^\circ$.
\item The following numbers are measures of angles in radians. What are their
      measures in degrees? $\pi/4$; $5\pi/4$; 3; 2; 1.5; $\pi/3$; $\pi/6$;
      $\pi/2$.
\item Which is the larger unit, radian or degree? For which unit will the
      measure of an angle give the larger number?
\item Consider a central angle of 1 radian on a circle of radius $r$. What is
      the length of the arc subtended by this angle? an angle of 2 radians? an
      angle of 3 radians? an angle of $n$ radians?
\end{exlist}
\end{multicols}

\par\medskip\noindent\emph{Consequence of Definition} 12--7.\ If $R$ is the
measure in radians of a central angle in a circle of radius $r$, then the angle
subtends an arc of length $r \cdot R$.\par\medskip

%========================================================= 12-9
\section*{12--9\quad OTHER UNITS}
\addcontentsline{toc}{section}{12--9\quad Other units}

Units other than degrees and radians could be chosen. One which has some use in
artillery is the \emph{mil}. It is approximately the angle which subtends an arc
of 1 yd at a distance of 1000 yd. Apart from its military use, it is of little
importance.

%========================================================= 12-10
\section*{12--10\quad INSCRIBED ANGLES}
\addcontentsline{toc}{section}{12--10\quad Inscribed angles}

Now that we have a measure for angles, it is easy to prove a theorem about the
measure of angles \emph{inscribed} in a circle. We recall the definition.

An angle is said to be \emph{inscribed} in a circle if its vertex and one point
on each side of the angle, other than the vertex, lie on the circle. The arc of
the circle interior to the angle is said to be subtended by the angle. (See
Fig.~12--19.)

\begin{theorem}
The measure of an inscribed angle in degrees is equal to one half of the degree
measure of the arc it subtends.
\end{theorem}

\FIGXIININETEENTWENTY

First, we prove a special case of this theorem, which will enable us to prove
the theorem easily.

\par\medskip\noindent\textbf{Lemma 12--11--1.}\ \emph{If $\ang{ABC}$ is an
inscribed angle and $BC$ is a diameter of the circle, then measure
$\ang{B} = \frac{1}{2}$ measure $\XIIarc{AC}$.} (See Fig.~12--20.)\par\medskip

\noindent\emph{Proof.}
\begin{steps}
\step{$\ang{OAB} \cong \ang{B}$.}{Why?}
\step{Measure $\ang{OAB} + $ measure $\ang{OBA} = $ measure $\ang{COA}$.}{Why?}
\step{2 (measure $\ang{B}$) $= $ measure $\ang{COA}$.}{Why?}
\step{Measure $\ang{COA} = $ measure $\XIIarc{AC}$.}{Why?}
\step{Measure $\ang{B} = \frac{1}{2}$ (measure $\XIIarc{AC}$).}{Why?}
\end{steps}

\noindent\emph{Proof of the theorem.} There are two cases to consider to
complete the proof of the theorem. They are: \emph{Case} 1. The center of the
circle is interior to $\ang{ABC}$ (Fig.~12--21). \emph{Case} 2. The center of
the circle is exterior to $\ang{ABC}$ (Fig.~12--22).

\FIGXIITWENTYONETWENTYTWO

In Case 1 the diameter through $B$ divides $\XIIarc{AC}$ into arcs
$\XIIarc{AP}$ and $\XIIarc{PC}$, and $\ang{ABC}$ into angles $x$ and $y$, as
shown. By Lemma 12--11--1, measure $\ang{x} = \frac{1}{2}$ (measure
$\XIIarc{AP}$), and measure $\ang{y} = \frac{1}{2}$ (measure $\XIIarc{PC}$).
Hence measure $\ang{ABC} = $ measure $\ang{x} + $ measure
$\ang{y} = \frac{1}{2}$ (measure $\XIIarc{AP}$) $+ \frac{1}{2}$ (measure
$\XIIarc{PC}$) $= \frac{1}{2}$ (measure $\XIIarc{AC}$).

The proof of Case 2 is left to the student.

\exercisegroup{12--11}
\begin{multicols}{2}
\begin{exlist}
\item If $\ang{ABC}$ is inscribed in a circle and measure
      $\XIIarc{AC} = 80^\circ$, how many degrees in $\ang{B}$?
\item An angle of $100^\circ$ is inscribed in a circle. How many degrees are
      there in the arc it subtends?
\item Prove that an angle ``inscribed in a semicircle'' is a right angle
      (Fig.~12--23).
\item How large can an angle be that is inscribed in a circle?
\item Prove that two inscribed angles in one circle which subtend arcs of equal
      measure are congruent.
\item Prove that a parallelogram inscribed in a circle must be a rectangle. More
      generally, prove that if a quadrilateral is inscribed in a circle, then the
      opposite angles are supplementary.
\item In Fig.~12--24 lines $AB$ and $CD$ are parallel. Prove that measure
      $\XIIarc{AC} = $ measure $\XIIarc{BD}$.
\item A trapezoid is inscribed in a circle. Prove that the base angles are
      congruent. [\emph{Hint:} Use the result of Exercise 7.]
\item[\stex 9.] In Fig.~12--25, $AC$ is tangent to the circle. Prove that
      measure $\ang{ABD} = \frac{1}{2}$ measure $\XIIarc{BD}$. [\emph{Hint:}
      Draw $DE \pll AC$. Show that arcs $\XIIarc{BD}$ and $\XIIarc{BE}$ have the
      same measure. Then use $\ang{ABD} \cong \ang{BDE} \cong \ang{BED}$.]
\item[\stex 10.] Prove, for Fig.~12--26, that measure
      $\ang{A} = \frac{1}{2}$ (measure $\XIIarc{CD} - $ measure $\XIIarc{BE}$).
      State this theorem in words.
\item[\stex 11.] State and prove theorems similar to that of Exercise 10 for
      cases in which either $AC$ or $AD$, or both, are tangent to the circle.
\item[\stex 12.] Prove, for Fig.~12--27, that measure
      $\ang{AEB} = \frac{1}{2}$ (measure $\XIIarc{AB} + $ measure
      $\XIIarc{DC}$). [\emph{Hint:} Draw $BC$.] State this theorem in words.
\item[13.] Using Fig.~12--27, prove that $\tri{AED} \simto \tri{BEC}$. Show that
      $\sg{AE} \cdot \sg{EC} = \sg{BE} \cdot \sg{ED}$. State this last property
      as a theorem.
\item[\stex 14.] Using Fig.~12--26, prove that $\tri{AEC} \simto \tri{ABD}$.
      Hence prove that $\sg{AB} \cdot \sg{AC} = \sg{AE} \cdot \sg{AD}$. State
      this last property as a theorem.
\item[15.] Two chords of a circle intersect as in Fig.~12--28. If
      $\sg{AO} = 12$, $\sg{OB} = 4$, and $\sg{CO} = 8$, find $\sg{OD}$.
      [\emph{Hint:} Use Exercise 13.]
\item[\stex 16.] From point $A$ outside a circle, secants are drawn as in
      Fig.~12--29. If $\sg{AD} = 6$, $\sg{DE} = 10$, and $\sg{AB} = 8$, find
      $\sg{BC}$.
\item[\stex 17.] In Fig.~12--30, $AB$ is tangent to the circle at $B$. Prove
      that $\tri{ACB} \simto \tri{ABD}$. You will need the result of Exercise 9
      for this. Prove also that $\sg{AB}^2 = \sg{AD} \cdot \sg{AC}$.
\item[18.] If, in Fig.~12--30, $\sg{AC} = 8$ and $\sg{CD} = 10$, determine
      $\sg{AB}$.
\end{exlist}
\end{multicols}

\FIGXIITWENTYTHREE

\FIGXIITWENTYFOURTWENTYFIVE

\FIGXIITWENTYSIXTWENTYSEVEN

\FIGXIITWENTYEIGHT

\FIGXIITWENTYNINETHIRTY

%========================================================= review
\par\bigskip
\begin{center}\bfseries\large REVIEW OF CHAPTER 12\end{center}
\addcontentsline{toc}{section}{Review of Chapter 12}
\nopagebreak

In this chapter definitions have been given of the following terms; be sure that
you can state them in your own words: arc, subtended arc, inscribed polygonal
path, arc length, degree, measure of angle in degrees, measure of arc in
degrees, radian, measure of angle in radians, inscribed angle.

The following is an outline of the important definitions and theorems of the
chapter. Study the outline so that you understand clearly how each step leads to
the next.

\par\smallskip\noindent\emph{Theorem:} The limit of the lengths of a sequence of
regular inscribed polygons exists.

\noindent\emph{Definition:} Arc length.

\noindent\emph{Theorem:} Can use inscribed polygons other than regular ones to
define arc length.

\noindent\emph{Theorem:} Length of ``sum'' of two arcs equals sum of lengths.

\noindent\emph{Theorem:} Greater angle subtends longer arc.

\noindent\emph{Theorem:} Angles congruent if and only if arc lengths equal.

\noindent\emph{Theorem:} Arc length is proportional to radius.

\noindent\emph{Theorem:} For every positive number less than the circumference,
there is an arc having length equal to the number.

\noindent\emph{Definition:} Measure of angles in degrees.

\noindent\emph{Definition:} Measure of angles in radians.

\noindent\emph{Theorems:} Angle measurement and counting.

\noindent\emph{Theorem:} Measure of inscribed angle is one half measure of
subtended arc.

\par\medskip
\begin{center}\bfseries Important Remark\end{center}
\nopagebreak

The purpose of the chapter has been to define measurement of angles. To make the
definition, we first defined \emph{arc length}, and then \emph{angle
measurement} in terms of an arc length.

\par\medskip
\begin{center}\bfseries Review Exercises\end{center}
\nopagebreak
\begin{exlist}
\item From the definition of arc length, show that the length of arc subtended
      by a right angle (Fig.~12--31) is greater than $\sqrt{2}\,r$. Hence,
      conclude that $\pi/2 > \sqrt{2}$.
\end{exlist}

\FIGXIITHIRTYONE

\begin{exlist}\setcounter{exlisti}{1}
\item Angles of $15^\circ$ subtend arcs of lengths 10 and 12 on two circles.
      What is the ratio of the radii of the circles? Find the radii.
\item Theorem 12--7 states that arc length is proportional to the radius. What
      theorem about triangles is used in the proof of Theorem 12--7?
\item State briefly what is involved in proving that the length of the ``sum''
      of two adjacent arcs is equal to the sum of their lengths.
\item Compute the length of an arc in a circle of radius 10 if it is subtended
      by an angle of (a) 40 degrees, (b) 2 radians, (c) $\pi/4$ radians,
      (d) $\frac{1}{2}$ degree.
\item[\stex 6.] A sequence of points $A_1$, $A_2$, $A_3$, \ldots\ is marked on a
      circle so that each pair of consecutive points determines a central angle
      of 1 rad. The direction from $A_1$ to $A_2$ to $A_3$, etc., is
      counterclockwise. Compute the approximate length of the arc from $A_{51}$
      to $A_1$ in a counterclockwise sense if the radius of the circle is 10.
      Use $\pi = 3.1416$.
\end{exlist}

\end{document}
